% concatenated from beta-28-avril.tex
%2multibyte Version: 5.50.0.2960 CodePage: 1250
\documentclass[reqno,openbib,runningheads,a4paper,12pt]{amsart}%
\usepackage{graphicx}
\usepackage{amssymb}
\usepackage{amsfonts}
\usepackage{amsmath}
\usepackage{graphicx}
\usepackage{lineno}
\usepackage[authoryear]{natbib}
\usepackage[dvipsnames]{xcolor}
\usepackage{epsf}
\usepackage{lscape}
\usepackage[legalpaper,bookmarks=true,colorlinks=true,linkcolor=blue,citecolor=blue]%
{hyperref}%
\setcounter{MaxMatrixCols}{30}
%TCIDATA{OutputFilter=latex2.dll}
%TCIDATA{Version=5.50.0.2960}
%TCIDATA{Codepage=1250}
%TCIDATA{LastRevised=Wednesday, May 13, 2026 15:48:35}
%TCIDATA{<META NAME="GraphicsSave" CONTENT="32">}
%TCIDATA{<META NAME="SaveForMode" CONTENT="1">}
%TCIDATA{BibliographyScheme=Manual}
%TCIDATA{Language=American English}
%TCIDATA{ComputeDefs=
%$P\left(  Z>x\right)  $
%}
%BeginMSIPreambleData
\providecommand{\U}[1]{\protect\rule{.1in}{.1in}}
%EndMSIPreambleData
\providecommand{\U}[1]{\protect \rule{.1in}{.1in}}
\textheight 248mm
\topmargin  0mm
\textwidth 165mm
\oddsidemargin 0mm
\evensidemargin 0mm
\def \theequation{\arabic{section}.\arabic{equation}}
\def \thetable{\arabic{section}.\arabic{table}}
\def \thefigure{\arabic{section}.\arabic{figure}}
\newtheorem{theorem}{Theorem}[section]
\newtheorem{corollary}{Corollary}[section]
\newtheorem{proposition}{Proposition}[section]
\newtheorem{statement}{Statement}[section]
\newtheorem{lemma}{Lemma}[section]
\newtheorem{example}{Example}[section]
\newtheorem{definition}{Definition}[section]
\newtheorem{note}{Note}[section]
\newtheorem{assumption}{Assumption}[section]
\newtheorem{remark}{Remark}[section]
\renewcommand{\baselinestretch}{1.4}
\makeatletter
\renewcommand{\@biblabel}[1]{}
\makeatother
\begin{document}

\begin{center}
{\large \textbf{Weighted and Truncated Tail Index Estimation under Random
Censoring: A Unified Full-Range Framework}}\bigskip

{\large Abdelhakim Necir}$^{\ast}${\large , Nour Elhouda Guesmia, Djamel
Meraghni}\medskip\newline

{\small \textit{Laboratory of Applied Mathematics, Mohamed Khider University,
Biskra, Algeria}}\medskip\medskip
\end{center}

\noindent\textbf{Abstract.} Estimation of the extreme value index under right
censoring is a fundamental problem in extreme value theory, with important
applications in finance, insurance, and reliability. Classical integral
estimators for Pareto-type tails typically require that the asymptotic
proportion $p$ of uncensored observations in the tail exceeds one half,
corresponding to the weak censoring regime. This restriction excludes many
practically relevant situations with strong censoring $(0<p\leq1/2)$ and
reflects the lack of a uniformly valid Gaussian approximation for the
associated tail empirical process. To overcome this limitation, we introduce a
weighted and truncated Nelson--Aalen tail empirical process and construct a
class of integral estimators indexed by a parameter $\beta>1$. This approach
restores a tractable asymptotic structure over the full censoring range
$0<p<1$. Under standard regular variation conditions, we establish a uniform
Gaussian approximation and derive consistency and asymptotic normality without
restrictions on $p$. A key ingredient is a linearization of the estimator as a
functional of the process. Simulation studies and real data applications
demonstrate improved stability and accuracy, particularly under moderate and
strong censoring. In particular, the analysis of insurance loss data (weak
censoring) and Australian AIDS survival data (strong censoring) highlights the
practical relevance of the proposed methodology across contrasting censoring
regimes.\medskip

\noindent\textbf{Keywords and Phrases:} Extreme value index; Random
Right-Censoring; Weak convergence.\medskip

\noindent\textbf{AMC 2020 Subject Classification:} 62G32; 62G05; 62G20.

\vfill


\vfill


\noindent{\small $^{\text{*}}$Corresponding author:
ah.necir@univ-biskra.dz\newline\noindent\textit{E-mail address:}\newline
nourelhouda.guesmia@univ-biskra.dz\texttt{\ }(N.~Guesmia)\newline
djamel.meraghni@univ-biskra.dz (D.~Meraghni)}\newline

\section{\textbf{Introduction\label{sec1}}}

\noindent Right-censored data arise naturally in many areas of applied
statistics, including actuarial science, survival analysis, reliability
engineering, and finance. In such settings, the variable of interest cannot
always be fully observed and is only known up to a censoring threshold. More
precisely, the true value of an observation is observed only if it occurs
before a censoring time and remains unknown otherwise. This partial
observation mechanism induces a non-negligible loss of information, thereby
requiring the development of dedicated statistical methods.\smallskip

\noindent In many practical applications, the underlying distributions exhibit
heavy tails. Pareto-type distributions are commonly used to model such
phenomena, as they allow for a significantly higher probability of extreme
observations compared with light-tailed models such as the normal
distribution. These distributions play a central role in the analysis of rare
and high-impact events and arise naturally in risk management, insurance,
finance, and environmental sciences, where extreme observations may have
substantial consequences.\smallskip

\noindent Let $X_{1},\ldots,X_{n}$ be a sample from a random variable $X$, and
$C_{1},\ldots,C_{n}$ a sample from another random variable $C$, defined on a
probability space $(\Omega,\mathcal{A},\mathbb{P})$, with continuous
cumulative distribution functions (cdfs) $F$ and $G$, respectively. Throughout
the paper, we assume that $X$ and $C$ are independent. We consider the
classical random right-censoring model where, for each $1\leq j\leq n$, we
observe the pair $(Z_{j},\delta_{j})$ defined by
\[
Z_{j}:=\min(X_{j},C_{j}), \qquad\delta_{j}:=\mathbb{I}_{\{X_{j}\leq C_{j}\}}.
\]
The indicator variable $\delta_{j}$ specifies whether the $j$-th observation
is uncensored. Hence, the observed sample combines exact observations and
censored values, leading to an incomplete description of the tail
behavior.\smallskip

\noindent We assume that the survival functions $\overline{F}:=1-F$ and
$\overline{G}:=1-G$ are regularly varying at infinity with negative indices
$-1/\gamma_{1}$ and $-1/\gamma_{2}$, respectively. Equivalently, for every
$x>0$,
\begin{equation}
\lim_{t\rightarrow\infty} \frac{\overline{F}(tx)}{\overline{F}(t)} =
x^{-1/\gamma_{1}}, \label{RVF}%
\end{equation}
and
\begin{equation}
\lim_{t\rightarrow\infty} \frac{\overline{G}(tx)}{\overline{G}(t)} =
x^{-1/\gamma_{2}}, \label{RVG}%
\end{equation}
These conditions imply that both $X$ and $C$ have Pareto-type tails, with tail
indices $\gamma_{1}$ and $\gamma_{2}$, respectively. In particular, the larger
the tail index, the heavier the corresponding tail, which directly influences
the frequency and magnitude of extreme observations.\smallskip

\noindent In order to derive refined asymptotic results, it is standard in
extreme value theory to assume a second-order framework; see, for instance,
\cite{deHF06}. Accordingly, we assume that $\overline{F}$ satisfies a
second-order condition of regular variation. This condition allows one to
quantify the rate of convergence of the Pareto approximation and plays a key
role in bias analysis. That is, for any $x>0$,
\[
\lim_{t\rightarrow\infty} \frac{U_{F}(tx)/U_{F}(t)-x^{\gamma_{1}}}{A_{1}%
^{\ast}(t)} = x^{\gamma_{1}}\frac{x^{\tau_{1}}-1}{\tau_{1}},
\]
or equivalently
\begin{equation}
\lim_{t\rightarrow\infty} \frac{\overline{F}(tx)/\overline{F}(t)-x^{-1/\gamma
_{1}}}{A_{1}(t)} = x^{-1/\gamma_{1}}\frac{x^{\tau_{1}/\gamma_{1}}-1}{\tau
_{1}\gamma_{1}}, \label{second-orderF}%
\end{equation}
where $\tau_{1}\leq0$ denotes the second-order parameter and $A_{1}(t)$ is a
function tending to zero at infinity. The parameter $\tau_{1}$ governs the
speed of convergence toward the Pareto limit.\smallskip\noindent The
functions
\[
L^{\leftarrow}(s):=\inf\{x:L(x)\geq s\},\qquad0<s<1,
\]
and
\[
U_{L}(t):=L^{\leftarrow}(1-1/t),\qquad t>1,
\]
denote, respectively, the quantile function and the tail quantile function
associated with a cdf $L$. The function $U_{L}$ provides a convenient
representation of the upper tail and plays a central role in extreme value
analysis.\smallskip

\noindent The second-order condition is satisfied by a wide class of
heavy-tailed distributions frequently used in practice, including Burr,
Fr\'{e}chet, generalized Pareto, generalized extreme value, and Student
distributions. This highlights the broad applicability of the theoretical
framework. However, the class of distributions satisfying
\eqref{second-orderF} is even broader. For instance, the modified Pareto
distribution
\[
F(x)=1-x^{-1/\gamma_{1}}\left(  1+\frac{1}{\log x}\right)  , \qquad x>e,
\]
satisfies \eqref{second-orderF} with $A_{1}(t)=1/\log t$. This example
illustrates the flexibility of the second-order condition beyond classical
parametric families.\smallskip

\noindent Let $Z=\min(X,C)$ and denote by $H$ its distribution function. Under
the independence assumption between $X$ and $C$, we have
\[
\overline{H}(x)=\overline{F}(x)\,\overline{G}(x),
\]
which implies that $\overline{H}$ is regularly varying with index $-1/\gamma$,
where
\[
\gamma:=\frac{\gamma_{1}\gamma_{2}}{\gamma_{1}+\gamma_{2}}.
\]
Thus, the observed variable $Z$ is also heavy-tailed, with tail index $\gamma
$. This harmonic-type relationship reflects the loss of tail information
induced by censoring and plays a central role in the asymptotic behavior of
estimators based on censored data.\smallskip

\noindent Estimation of the tail index in the presence of right-censored data
has received considerable attention in the literature. Several authors have
proposed adaptations of classical extreme value estimators to accommodate
censoring. Most of these approaches rely on modifying standard estimators by
incorporating information about the censoring mechanism. A well-known example
is the estimator introduced by \cite{EnFG08}, which adapts the Hill estimator
\citep{Hill75} to the censoring framework. It is defined by
\[
\widehat{\gamma}_{1,k}^{(\mathrm{EFG})}:=\frac{\widehat{\gamma}_{k}%
^{(\mathrm{H})}}{\widehat{p}_{k}},
\]
where
\[
\widehat{\gamma}_{k}^{(\mathrm{H})}:=\frac{1}{k}\sum_{i=1}^{k}\log
\frac{Z_{n-i+1:n}}{Z_{n-k:n}}%
\]
is the classical Hill estimator based on the observed sample $Z$, and
\[
\widehat{p}_{k}:=\frac{1}{k}\sum_{i=1}^{k}\delta_{\lbrack n-i+1:n]}%
\]
estimates the proportion of uncensored observations in the tail,
\[
p:=\frac{\gamma_{2}}{\gamma_{1}+\gamma_{2}}=\frac{\gamma}{\gamma_{1}}.
\]
The parameter $p$ represents the asymptotic proportion of uncensored
observations in the upper tail and plays a central role in the behavior of
tail index estimators under censoring. In particular, the regime $p>1/2$
corresponds to weak censoring, whereas $p\leq1/2$ characterizes moderate to
strong censoring. In the latter case, the effective amount of tail information
becomes too limited, and classical estimators may suffer from instability or
even asymptotic breakdown.\smallskip

\noindent Here $k=k_{n}$ denotes the number of upper order statistics,
satisfying $k\rightarrow\infty$ and $k/n\rightarrow0$ as $n\rightarrow\infty$.
The order statistics of the sample $Z_{1},\ldots,Z_{n}$ are denoted by
$Z_{1:n}\leq\cdots\leq Z_{n:n}$, and $\delta_{\lbrack1:n]},\ldots
,\delta_{\lbrack n:n]}$ represent the associated concomitants. More precisely,
$\delta_{\lbrack i:n]}$ is the censoring indicator corresponding to the
observation $Z_{i:n}$, that is, the value of $\delta_{j}$ such that
$Z_{j}=Z_{i:n}$.\smallskip

\noindent This intermediate sequence framework is standard in extreme value
theory and ensures an appropriate balance between bias and variance in tail
estimation.\smallskip

\noindent Gaussian approximations for $\widehat{p}_{k}$, $\widehat{\gamma}%
_{k}^{(\mathrm{H})}$, and $\widehat{\gamma}_{1,k}^{(\mathrm{EFG})}$ in terms
of Brownian bridges were obtained in \cite{BMN-2015}. These results provide a
refined description of the stochastic fluctuations of the estimators and form
the basis for deriving asymptotic normality.\smallskip

\noindent Another important line of research relies on Kaplan--Meier
integration. This approach exploits the product-limit structure of the
Kaplan--Meier estimator to handle censoring in a fully nonparametric way. In
this direction, \cite{WW2014} proposed the estimator
\[
\widehat{\gamma}_{1,k}^{(\mathrm{W})}:=\sum_{i=1}^{k}\frac{\delta_{\lbrack
n-i+1:n]}}{i}\frac{\overline{F}_{n}^{(\mathrm{KM})}(Z_{n-i+1:n})}{\overline
{F}_{n}^{(\mathrm{KM})}(Z_{n-k:n})}\log\frac{Z_{n-i+1:n}}{Z_{n-k:n}},
\]
where
\[
F_{n}^{(\mathrm{KM})}(z):=1-\prod_{Z_{i:n}\leq z}\left(  1-\frac
{\delta_{\lbrack i:n]}}{n-i+1}\right)
\]
denotes the Kaplan--Meier estimator \citep{KM58}.This estimator can be
interpreted as a weighted version of the Hill estimator, where the weights are
induced by the Kaplan--Meier estimate of the survival function, thereby
correcting for the censoring mechanism in the upper tail.\smallskip

\noindent Recently, \cite{BR2025} introduced an estimator based on an extreme
Kaplan--Meier construction, which can be written explicitly as
\[
\widehat{\gamma}_{1,k}^{(\mathrm{BR})}:=\sum_{i=1}^{k}\frac{\delta_{\lbrack
n-i+1:n]}}{i}\prod_{j=i+1}^{k}\left(  \frac{j-1}{j}\right)  ^{\delta_{\lbrack
n-j+1:n]}}\log\frac{Z_{n-i+1:n}}{Z_{n-k:n}}.
\]
At first glance, this representation appears structurally different from the
Kaplan--Meier integral form, due to the presence of the explicit product term.
To clarify its relationship with the estimator of \cite{WW2014}, consider the
multiplicative weight appearing in $\widehat{\gamma}_{1,k}^{(\mathrm{BR})}$:
\[
\prod_{j=i+1}^{k}\left(  \frac{j-1}{j}\right)  ^{\delta_{\lbrack n-j+1:n]}%
}=\prod_{m=n-k}^{n-i-1}\left(  1-\frac{1}{m+1}\right)  ^{\delta_{\lbrack
m:n]}}.
\]
This reindexing expresses the product in terms of the original order
statistics, making its probabilistic interpretation more transparent. The
change of index $m=n-j$ shows that this product coincides with the
product-limit structure underlying the Kaplan--Meier estimator. More
precisely, it corresponds to the multiplicative weights defining $\overline
{F}_{n}^{(\mathrm{KM})}$ evaluated over the upper order statistics. In other
words, the product term represents the discrete analogue of the Kaplan--Meier
survival function restricted to the tail region. Consequently,
\[
\widehat{\gamma}_{1,k}^{(\mathrm{BR})}\equiv\widehat{\gamma}_{1,k}%
^{(\mathrm{W})},
\]
that is, the estimator of \cite{BR2025} coincides with the Kaplan--Meier
integral estimator introduced by \cite{WW2014}. This equivalence shows that
both estimators rely on the same underlying mechanism, despite their different
algebraic representations.\smallskip

\noindent More recently, \cite{MNS2025} introduced a tail index estimator
based on a Nelson--Aalen integral representation. The main motivation behind
this approach is to obtain a formulation leading to a simpler and more
transparent asymptotic analysis while preserving the main characteristics of
Kaplan--Meier based procedures. The resulting estimator is defined by
\[
\widehat{\gamma}_{1,k}^{(\mathrm{MNS})}:=\sum_{i=1}^{k}\frac{\delta_{\lbrack
n-i+1:n]}}{i}\frac{\overline{F}_{n}^{(\mathrm{NA})}(Z_{n-i+1:n})}{\overline
{F}_{n}^{(\mathrm{NA})}(Z_{n-k:n})}\log\frac{Z_{n-i+1:n}}{Z_{n-k:n}},
\]
where
\[
F_{n}^{(\mathrm{NA})}(z):=1-\prod_{Z_{i:n}\leq z}\exp\!\left(  -\frac
{\delta_{\lbrack i:n]}}{n-i+1}\right)
\]
denotes the Nelson--Aalen estimator of the distribution function $F$
\citep{Nelson1972}.\smallskip

\noindent Overall, these results indicate that the choice between
Kaplan--Meier and Nelson--Aalen formulations is primarily driven by analytical
convenience rather than differences in asymptotic efficiency.\smallskip

\noindent The Nelson--Aalen formulation naturally induces a tail empirical
process that plays a role analogous to the extreme Kaplan--Meier process. In
\cite{MNS2025}, the authors derive an explicit Gaussian approximation for this
process, which yields a stochastic expansion of the estimator and allows them
to establish, in a unified and rigorous manner, both its consistency and
asymptotic normality. This representation separates the leading stochastic
fluctuations from higher-order bias terms, which is essential for asymptotic
analysis.\smallskip

\noindent More generally, this functional approach significantly simplifies
the derivation of asymptotic properties for extreme value statistics under
censoring. By working at the level of stochastic processes, it avoids the
technically intricate combinatorial arguments inherent in product-limit
representations, as encountered in \cite{WW2014} and \cite{BWW2019}.\smallskip

\noindent Despite their structural differences, the Kaplan--Meier-- and
Nelson--Aalen--based estimators share identical asymptotic bias and variance
under the standard weak censoring conditions. This theoretical equivalence is
further supported by empirical studies, which consistently report comparable
performance in terms of bias and mean squared error (MSE); see, for instance,
\cite{Colosimo2002}.\smallskip

\noindent More generally, this functional approach significantly simplifies
the derivation of asymptotic properties for extreme value statistics under
censoring. By working at the level of stochastic processes, it avoids the
technically intricate combinatorial arguments inherent in product-limit
representations, as encountered in \cite{WW2014} and \cite{BWW2019}.\smallskip

\noindent Both \cite{BR2025} and \cite{MNS2025} pursue the same fundamental
objective, namely the development of functional approaches to censored tail
index estimation that enable a rigorous yet simplified asymptotic analysis.
While \cite{BR2025} relies on the extreme Kaplan--Meier process,
\cite{MNS2025} exploits the Nelson--Aalen integral. Despite these different
constructions, both approaches lead to equivalent estimators and similar
asymptotic behaviors, reinforcing the idea that the choice between them is
mainly guided by analytical convenience.

\subsection{\textbf{Main contribution: weighted and truncated Nelson--Aalen
estimator}}

\noindent A major limitation of existing approaches is that their theoretical
validity relies on the condition $p>1/2$, corresponding to the weak censoring
regime. This restriction is not merely technical but intrinsic to the Gaussian
approximation of the associated tail empirical processes. In contrast, many
practical situations involve moderate or strong censoring, that is,
$0<p\leq1/2$. In such cases, the Gaussian approximation is no longer uniformly
valid, which prevents the derivation of standard asymptotic results. As a
consequence, existing estimators may exhibit instability, inflated variance,
or even a breakdown of their asymptotic properties.\smallskip\noindent This
limitation motivates the development of a new methodology capable of ensuring
uniform asymptotic validity over the entire censoring range $0<p<1$. The key
idea of the present work is to restore a tractable stochastic structure for
the tail empirical process under arbitrary censoring levels by acting directly
at the level of the process itself rather than only at the level of the
estimator.\smallskip

\noindent The construction of the proposed estimator is rooted in the
classical identity (see \cite{MNS2025})
\[
\int_{1}^{\infty}x^{-1}\frac{\overline{F}(tx)}{\overline{F}(t)}%
\,dx\longrightarrow\gamma_{1},\qquad t\rightarrow\infty,
\]
which provides a direct link between the tail index and an integral functional
of the normalized survival function. An equivalent representation, obtained by
integration by parts, is
\begin{equation}
\int_{t}^{\infty}\log\!\left(  \frac{x}{t}\right)  \,d\!\left(  \frac
{F(x)}{\overline{F}(t)}\right)  \longrightarrow\gamma_{1},\qquad
t\rightarrow\infty. \label{log}%
\end{equation}
Replacing $F$ by its Nelson--Aalen estimator and taking $t=Z_{n-k:n}$ yields
\begin{equation}
\widehat{\gamma}_{1,k}^{(\mathrm{MNS})}=\int_{Z_{n-k:n}}^{\infty}\log\!\left(
\frac{x}{Z_{n-k:n}}\right)  \,d\!\left(  \frac{F_{n}^{(\mathrm{NA})}%
(x)}{\overline{F}_{n}^{(\mathrm{NA})}(Z_{n-k:n})}\right)  . \label{gchap}%
\end{equation}
After a change of variables, this estimator admits the functional
representation
\[
\widehat{\gamma}_{1,k}^{(\mathrm{MNS})}=\int_{1}^{\infty}x^{-1}\frac
{\overline{F}_{n}^{(\mathrm{NA})}(xZ_{n-k:n})}{\overline{F}_{n}^{(\mathrm{NA}%
)}(Z_{n-k:n})}\,dx.
\]
Introducing the Nelson--Aalen tail process
\[
D_{n}(x):=\sqrt{k}\left\{  \frac{\overline{F}_{n}^{(\mathrm{NA})}(xZ_{n-k:n}%
)}{\overline{F}_{n}^{(\mathrm{NA})}(Z_{n-k:n})}-x^{-1/\gamma_{1}}\right\}  ,
\]
we obtain
\[
\sqrt{k}\left(  \widehat{\gamma}_{1,k}^{(\mathrm{MNS})}-\gamma_{1}\right)
=\int_{1}^{\infty}x^{-1}D_{n}(x)\,dx.
\]
This representation shows that the estimator is a linear functional of the
tail process. Such a structure is crucial for asymptotic analysis, as it
reduces the problem to the weak convergence of the underlying stochastic
process. However, this approach critically relies on the uniform validity of
the Gaussian approximation of $D_{n}(x)$, which typically holds only when
$p>1/2$.\smallskip

\noindent When $p\leq1/2$, this approximation breaks down due to two distinct phenomena:

\begin{itemize}
\item the remainder terms are no longer uniformly negligible,

\item a pathological contribution arises from the lower part of the
transformed scale.
\end{itemize}

\noindent These effects prevent a stable normalization of the estimator and
invalidate the classical asymptotic framework. Intuitively, under strong
censoring, the effective amount of uncensored information in the tail becomes
too limited to ensure a regular Gaussian behavior.\smallskip

\noindent To overcome this limitation, we introduce a weighted Nelson--Aalen
tail process defined on a truncated domain:
\[
D_{n,a}^{(\beta)}(x)=\sqrt{k}\,x^{-\beta/\gamma}\left\{  \frac{\overline
{F}_{n}^{(\mathrm{NA})}(xZ_{n-k:n})}{\overline{F}_{n}^{(\mathrm{NA}%
)}(Z_{n-k:n})}-x^{-1/\gamma_{1}}\right\}  ,\qquad1\leq x\leq T_{n,a},
\]
where
\[
u_{n}:=Q^{(1)}(a_{n}),\qquad T_{n,a}:=\frac{u_{n}}{Z_{n-k:n}},
\]
and $(a_{n})$ satisfies
\begin{equation}
a_{n}\rightarrow0,\qquad na_{n}\rightarrow\infty,\qquad a_{n}=o(k/n).
\label{assump_an}%
\end{equation}
Here $Q^{(1)}$ denotes the (generalized) quantile function associated with the
sub-distribution
\[
H^{(1)}(z):=\mathbb{P}(Z\leq z,\ \delta=1),
\]
and is defined by
\[
Q^{(1)}(s):=\inf\left\{  z:\ H^{(1)}(z)\geq s\right\}  ,\qquad0<s<H^{(1)}%
(\infty).
\]
\noindent We also introduce the complementary sub-distribution corresponding
to censored observations, defined by
\[
H^{(0)}(z):=\mathbb{P}(Z\leq z,\ \delta=0).
\]
These two sub-distributions satisfy the fundamental decomposition
\[
H(z)=H^{(0)}(z)+H^{(1)}(z).
\]
The above conditions ensure that the truncation level $u_{n}$ diverges to
infinity while remaining asymptotically negligible relative to the
intermediate threshold $Z_{n-k:n}$. In particular, the condition
$na_{n}\rightarrow\infty$ guarantees that the truncated region contains enough
observations for a stable empirical approximation, whereas $a_{n}=o(k/n)$
ensures that the truncation does not affect the leading asymptotic behavior of
the estimator.\smallskip

\noindent The truncation is introduced at the level of the domain of the
process rather than through an explicit modification of the numerator. This is
a crucial point: it removes the pathological contribution arising from the
lower part of the transformed scale without generating additional boundary
terms. The weight $x^{-\beta/\gamma}$ further controls the stochastic
fluctuations and improves the decay of the remainder terms.\smallskip

\noindent The main contribution of this paper is to restore a valid asymptotic
framework for tail index estimation under random censoring over the full range
$0<p<1$, including the strong censoring regime where classical methods
fail.\smallskip

\noindent\ To achieve this, we introduce a weighted and truncated
Nelson--Aalen approach that modifies the tail empirical process at a
structural level, thereby ensuring uniform asymptotic validity across all
censoring regimes.\smallskip

\noindent The proposed estimator is constructed as the empirical counterpart
of the weighted and truncated identity
\begin{equation}
\left(  \frac{\beta}{p}\right)  ^{2}\int_{t}^{u_{t}}\left(  \frac{\overline
{F}(y)}{\overline{F}(t)}\right)  ^{\beta/p-1}\log\!\left(  \frac{y}{t}\right)
\,d\!\left(  \frac{F(y)}{\overline{F}(t)}\right)  \longrightarrow\gamma_{1},
\label{ass2-tr}%
\end{equation}
which holds for every $\beta>1$ and $0<p<1$. See
Proposition~\ref{prop:app-identity}.\smallskip\noindent This construction
directly mirrors the continuous identity \eqref{ass2-tr}, ensuring that the
estimator arises as its natural empirical counterpart rather than being
introduced in an ad hoc manner. More precisely, the empirical version is
obtained by replacing $t$ with $Z_{n-k:n}$, $u_{t}$ with $u_{n}$, and $F$ with
its Nelson--Aalen estimator $F_{n}^{(\mathrm{NA})}$. The resulting discrete
form is given by
\[
\widehat{\gamma}_{1,k}^{(\mathrm{NA,tr})}(\beta)=\left(  \frac{\beta
}{\widehat{p}_{k}}\right)  ^{2}\sum_{i=m_{n}}^{k}\frac{\delta_{\lbrack
n-i+1:n]}}{i}\prod_{j=i+1}^{k}\exp\!\left\{  \left(  1-\frac{\beta
}{\widehat{p}_{k}}\right)  \frac{\delta_{\lbrack n-j+1:n]}}{j}\right\}
\log\!\left(  \frac{Z_{n-i+1:n}}{Z_{n-k:n}}\right)  ,
\]
where $m_{n}$ is the truncation index defined by $Z_{n-m_{n}+1:n}\leq
u_{n}<Z_{n-m_{n}:n}$. A natural and convenient choice is
\[
m_{n}:=\lfloor na_{n}\rfloor.
\]
Under these conditions, we have
\[
m_{n}\rightarrow\infty,\qquad\frac{m_{n}}{k}\rightarrow0,
\]
so that the truncation removes an asymptotically diverging but negligible
fraction of upper order statistics. Moreover, this choice is consistent with
the continuous truncation level $u_{n}$, since $Z_{n-m_{n}+1:n}\approx u_{n}%
$.\smallskip

\noindent For more details on the construction of the estimator and the
practical choice of $m_{n}$, see Subsections \ref{construct} and
\ref{choice_mn}, respectively.\smallskip

\noindent Hence, the proposed estimator is the exact empirical counterpart of
the weighted and truncated identity \eqref{ass2-tr}, providing a principled
extension of classical Nelson--Aalen-based tail index estimators to the full
censoring range $0<p<1$.\smallskip

\noindent A key feature of the proposed approach is that the estimator admits
the linearization
\begin{equation}
\sqrt{k}\left(  \widehat{\gamma}_{1,k}^{(\mathrm{NA,tr})}(\beta)-\gamma
_{1}\right)  =\left(  \frac{\beta}{p}\right)  ^{2}\int_{1}^{T_{n,a}%
}x^{-1+1/\gamma_{1}}D_{n,a}^{(\beta)}(x)\,dx+B_{n}+o_{\mathbf{P}}(1),
\label{functional-key}%
\end{equation}
where $B_{n}$ is a deterministic bias term.\smallskip\noindent This
representation is the cornerstone of the analysis. It shows that the estimator
is asymptotically equivalent to a linear functional of the weighted tail
process over the truncated domain, thereby restoring a Gaussian-type
asymptotic structure over the entire censoring range $0<p<1$. In particular,
it allows the Gaussian approximation of $D_{n,a}^{(\beta)}(x)$ to be directly
transferred to the estimator itself.\smallskip

\noindent Therefore, the joint use of domain truncation and weighting provides
a minimal yet effective modification of the classical framework, sufficient to
overcome the intrinsic limitations of existing approaches and to establish a
unified asymptotic theory under arbitrary censoring levels.

\subsubsection{\textbf{Constructing the estimator} $\protect\widehat{\gamma
}_{1,k}^{(\mathrm{NA,tr})}(\beta)$\textbf{\label{construct}}}

\noindent We now express the previous integral in a more tractable discrete
form based on the observed sample. To this end, we rely on the classical
representation of $F$ in terms of the distribution $H$ of $Z=\min(X,C)$ and
the sub-distribution $H^{(1)}$ associated with uncensored observations:
\[
\int_{0}^{z}\frac{dH^{(1)}(y)}{\overline{H}(y-)}=\int_{0}^{z}\frac
{dF(y)}{\overline{F}(y)}=-\log\overline{F}(z),\qquad z>0.
\]
Consequently,
\[
\overline{F}(z)=\exp\left\{  -\int_{0}^{z}\frac{dH^{(1)}(y)}{\overline{H}%
(y-)}\right\}  .
\]
The empirical counterparts are
\[
H_{n}(z):=\frac{1}{n}\sum_{i=1}^{n}\mathbb{I}_{\{Z_{i:n}\leq z\}},\qquad
H_{n}^{(1)}(z):=\frac{1}{n}\sum_{i=1}^{n}\delta_{\lbrack i:n]}\mathbb{I}%
_{\{Z_{i:n}\leq z\}},
\]
where $\mathbb{I}$ denotes the indicator function. It is worth noting that the
empirical distribution function $H_{n}$ can be decomposed into the two
sub-distribution functions $H_{n}^{(0)}$ and $H_{n}^{(1)}$ through the
identity
\[
H_{n}=H_{n}^{(0)}+H_{n}^{(1)}.
\]
Here
\[
H_{n}^{(0)}(z):=\frac{1}{n}\sum_{i=1}^{n}\left(  1-\delta_{\lbrack
i:n]}\right)  \mathbb{I}_{\{Z_{i:n}\leq z\}}.
\]
The function $H_{n}^{(0)}$ represents the empirical sub-distribution
associated with censored observations, while $H_{n}^{(1)}$ corresponds to
uncensored observations. These sub-distributions play a crucial role in the
theoretical analysis of random censoring, as they allow one to separate the
respective contributions of censored and observed data. For further details,
we refer to \cite{SW86}, p.~295. Finally, the Nelson--Aalen-type estimator of
$F$ is defined by
\begin{equation}
\overline{F}_{n}^{(\mathrm{NA})}(z):=\exp\left\{  -\int_{0}^{z}\frac
{dH_{n}^{(1)}(y)}{\overline{H}_{n}(y-)}\right\}  . \label{emp-F-formula-adapt}%
\end{equation}
\noindent Differentiating yields
\[
\frac{dF_{n}^{(\mathrm{NA})}(z)}{\overline{F}_{n}^{(\mathrm{NA})}(z)}%
=\frac{dH_{n}^{(1)}(z)}{\overline{H}_{n}(z-)}.
\]
Therefore,
\[
\widehat{\gamma}_{1,k}^{(\mathrm{NA,tr})}(\beta)=\left(  \frac{\beta
}{\widehat{p}_{k}}\right)  ^{2}\int_{0}^{\infty}\mathbb{I}_{\{Z_{n-k:n}%
<y<u_{n}\}}\left(  \frac{\overline{F}_{n}^{(\mathrm{NA})}(y)}{\overline{F}%
_{n}^{(\mathrm{NA})}(Z_{n-k:n})}\right)  ^{\beta/\widehat{p}_{k}-1}%
\log\!\left(  \frac{y}{Z_{n-k:n}}\right)  \frac{1}{\overline{H}_{n}%
(y-)}\,dH_{n}^{(1)}(y).
\]
Let $g(y)$ denote the integrand. Then, using the definition of $H_{n}^{(1)}$,
we obtain
\[
\int_{0}^{\infty}g(y)\,dH_{n}^{(1)}(y)=\frac{1}{n}\sum_{i=1}^{n}%
\delta_{\lbrack i:n]}\,g(Z_{i:n}).
\]
Moreover, recalling that $n\,\overline{H}_{n}(Z_{i:n}-)=n-i+1$, we deduce
that
\[
\widehat{\gamma}_{1,k}^{(\mathrm{NA,tr})}(\beta)=\left(  \frac{\beta
}{\widehat{p}_{k}}\right)  ^{2}\sum_{i=1}^{n}\mathbb{I}_{\{Z_{n-k:n}%
<Z_{i:n}<u_{n}\}}\left(  \frac{\overline{F}_{n}^{(\mathrm{NA})}(Z_{i:n}%
)}{\overline{F}_{n}^{(\mathrm{NA})}(Z_{n-k:n})}\right)  ^{\beta/\widehat{p}%
_{k}-1}\log\!\left(  \frac{Z_{i:n}}{Z_{n-k:n}}\right)  \frac{\delta_{\lbrack
i:n]}}{n-i+1}.
\]
The summation index arises from the restriction of the integration domain to
$[Z_{n-k:n},u_{n}]$, which corresponds exactly to the order statistics
satisfying $Z_{n-k:n}\leq Z_{i:n}\leq u_{n}$. \noindent Under the reversed
indexing, we write $Z_{i:n}=Z_{n-r+1:n}$, where $r=n-i+1$. Hence
\[
n-i+1=r,\qquad\delta_{\lbrack i:n]}=\delta_{\lbrack n-r+1:n]}.
\]
The restriction $Z_{n-k:n}<Z_{i:n}<u_{n}$ is therefore equivalent to
$r=m_{n},\ldots,k$. Consequently, the previous expression becomes
\[
\widehat{\gamma}_{1,k}^{(\mathrm{NA,tr})}(\beta)=\left(  \frac{\beta
}{\widehat{p}_{k}}\right)  ^{2}\sum_{r=m_{n}}^{k}\frac{\delta_{\lbrack
n-r+1:n]}}{r}\left(  \frac{\overline{F}_{n}^{(\mathrm{NA})}(Z_{n-r+1:n}%
)}{\overline{F}_{n}^{(\mathrm{NA})}(Z_{n-k:n})}\right)  ^{\beta/\widehat{p}%
_{k}-1}\log\!\left(  \frac{Z_{n-r+1:n}}{Z_{n-k:n}}\right)  .
\]
Renaming the index $r$ as $i$, we obtain
\[
\widehat{\gamma}_{1,k}^{(\mathrm{NA,tr})}(\beta)=\left(  \frac{\beta
}{\widehat{p}_{k}}\right)  ^{2}\sum_{i=m_{n}}^{k}\frac{\delta_{\lbrack
n-i+1:n]}}{i}\left(  \frac{\overline{F}_{n}^{(\mathrm{NA})}(Z_{n-i+1:n}%
)}{\overline{F}_{n}^{(\mathrm{NA})}(Z_{n-k:n})}\right)  ^{\beta/\widehat{p}%
_{k}-1}\log\!\left(  \frac{Z_{n-i+1:n}}{Z_{n-k:n}}\right)  .
\]
We now simplify the Nelson--Aalen survival ratio. From
\eqref{emp-F-formula-adapt}, for upper order statistics we have
\[
\frac{\overline{F}_{n}^{(\mathrm{NA})}(Z_{n-i+1:n})}{\overline{F}%
_{n}^{(\mathrm{NA})}(Z_{n-k:n})}=\prod_{j=i+1}^{k}\exp\!\left(  -\frac
{\delta_{\lbrack n-j+1:n]}}{j}\right)  .
\]
Therefore,
\[
\left(  \frac{\overline{F}_{n}^{(\mathrm{NA})}(Z_{n-i+1:n})}{\overline{F}%
_{n}^{(\mathrm{NA})}(Z_{n-k:n})}\right)  ^{\beta/\widehat{p}_{k}-1}%
=\prod_{j=i+1}^{k}\exp\!\left[  \left(  1-\frac{\beta}{\widehat{p}_{k}%
}\right)  \frac{\delta_{\lbrack n-j+1:n]}}{j}\right]  .
\]
Substituting this identity into the previous expression yields the fully
explicit discrete form
\[
\widehat{\gamma}_{1,k}^{(\mathrm{NA,tr})}(\beta)=\left(  \frac{\beta
}{\widehat{p}_{k}}\right)  ^{2}\sum_{i=m_{n}}^{k}\frac{\delta_{\lbrack
n-i+1:n]}}{i}\prod_{j=i+1}^{k}\exp\!\left[  \left(  1-\frac{\beta}%
{\widehat{p}_{k}}\right)  \frac{\delta_{\lbrack n-j+1:n]}}{j}\right]
\log\!\left(  \frac{Z_{n-i+1:n}}{Z_{n-k:n}}\right)  .
\]
As usual, an empty product is interpreted as equal to one.

\begin{remark}
A key feature of the proposed approach is that the modification is carried out
at the level of the tail empirical process rather than at the level of the
estimator itself. In contrast to existing methods, which adapt classical
estimators to the censored setting, we directly modify the stochastic
structure of the Nelson--Aalen tail process through a combination of
truncation and weighting. This perspective provides a unified framework in
which the estimator inherits its asymptotic properties from a suitably
regularized process. In particular, the functional representation
\eqref{functional-key} shows that the proposed estimator remains
asymptotically equivalent to a linear functional of a tail empirical process,
as in the classical uncensored case, but with a process specifically designed
to ensure a uniform Gaussian approximation over the full censoring range
$0<p<1$. Consequently, the present methodology should not be viewed as a
competing estimator but rather as a structural extension of the classical
Nelson--Aalen framework, restoring its asymptotic validity under strong
censoring through minimal and interpretable modifications.
\end{remark}

\subsubsection{\textbf{Practical choice of the truncation level }$m_{n}$
\textbf{\label{choice_mn}}}

\noindent The truncation index $m_{n}$ plays a crucial role in the
construction of the proposed estimator, as it removes the pathological
contribution arising from the lower part of the transformed scale. This region
corresponds to values for which the Gaussian approximation of the tail process
is no longer reliable.\smallskip

\noindent From a theoretical perspective, it is sufficient to assume that
$m_{n}\rightarrow\infty$ and $m_{n}/k\rightarrow0$, ensuring that the
truncation is asymptotically negligible while eliminating the unstable region.
These conditions guarantee that the truncation does not affect the first-order
asymptotic behavior of the estimator.\smallskip

\noindent In practice, the choice of $m_{n}$ must balance two competing
requirements. On the one hand, $m_{n}$ should diverge in order to remove the
contribution of the region near zero, where the Gaussian approximation fails
and induces a non-negligible bias under strong censoring. On the other hand,
$m_{n}$ must remain negligible compared to $k$, so that no relevant tail
information is lost.\smallskip

\noindent A convenient choice satisfying these requirements is
\[
m_{n}=\max\left(  3,\left\lfloor \log\log k\right\rfloor \right)  ,
\]
which diverges extremely slowly while ensuring that only a negligible fraction
of upper order statistics is discarded.\smallskip

\noindent An alternative choice is given by
\[
m_{n}=\left\lfloor k^{\rho}\right\rfloor ,\qquad0<\rho<1-\frac{1}{2\beta},
\]
but this approach requires the selection of an additional tuning parameter
$\rho$, which may complicate the practical implementation. In contrast, the
previous slowly varying choice is fully automatic, asymptotically valid, and
practically negligible.\smallskip

\noindent For instance, even for very large values of $k$, the quantity
$\log\log k$ remains small, so that $m_{n}$ typically takes values between $2$
and $4$ in practice. For example, for $n=2000$ (so that $k\approx200$), we
obtain $m_{n}=3$, showing that only a negligible number of upper order
statistics is discarded. This illustrates that the truncation has a negligible
impact on the empirical performance of the estimator.\smallskip

\noindent The truncation index is directly linked to the threshold $u_{n}$
through the relation $Z_{n-m_{n}+1:n}\approx u_{n}$, ensuring consistency
between the continuous and discrete formulations.\smallskip

\noindent The remainder of the paper is organized as follows.
Section~\ref{sec2} presents the main theoretical results, including the
Gaussian approximation of the weighted and truncated tail empirical process,
as well as the consistency and asymptotic normality of the proposed estimator.
This is followed by a simulation study in Section~\ref{sec3}, which evaluates
the performance of the estimator. Section~\ref{sec4} illustrates the
methodology using two real datasets: insurance loss data (weak censoring) and
AIDS survival times (strong censoring). These applications highlight the
reliability and robustness of the estimator across different tail-censoring
proportions and contamination scenarios. Proofs are deferred to
Section~\ref{sec5}, and Section~\ref{sec6} concludes the paper. Additional
technical results and simulation figures are provided in Appendices~A and~B, respectively.

\section{\textbf{Main results}\label{sec2}}

\noindent The following result is inspired by part of the proof of Theorem~2.1
in \cite{MNS2025}, but is adapted to the weighted and truncated Nelson--Aalen
framework introduced in the present paper. In contrast to the classical
setting, the present formulation explicitly incorporates both a reweighting
mechanism and a truncation device, which fundamentally modify the stochastic
structure of the tail process. Unlike the original result, which holds only
under the weak censoring condition $p>1/2$, the present formulation remains
valid over the entire censoring range $0<p<1$ by combining suitable weighting
with explicit truncation of the lower unstable tail region. This extension
constitutes a key methodological advance of the paper.\smallskip

\noindent The next three results constitute the main theoretical contribution
of the paper. They provide a complete and unified asymptotic characterization
of the proposed estimator, including a Gaussian approximation of the
underlying process, consistency, and asymptotic normality.

\begin{theorem}
\label{theorem1}

Assume that the survival function $\overline{F}$ satisfies the first-order
regular variation condition \eqref{RVF} and the second-order condition
\eqref{second-orderF} with auxiliary function $A_{1}$ such that
\[
\sqrt{k}\,A_{1}(h)\longrightarrow\lambda,\qquad h:=U_{H}(n/k).
\]
Let $k=k_{n}$ be such that $k\rightarrow\infty$ and $k/n\rightarrow0$, and let
$(a_{n})$ be a positive sequence satisfying
\[
a_{n}\rightarrow0,\qquad na_{n}\rightarrow\infty,\qquad a_{n}=o(k/n).
\]
Set
\[
Z_{k}:=Z_{n-k:n},\qquad T_{n,a}:=\frac{u_{n}}{Z_{k}},
\]
and let $\beta>1$. Define the weighted and truncated Nelson--Aalen tail
process by
\[
D_{n,a}^{(\beta)}(x):=\sqrt{k}\,x^{-\beta/\gamma}\left\{  \frac{\overline
{F}_{n}^{(\mathrm{NA})}(xZ_{k})}{\overline{F}_{n}^{(\mathrm{NA})}(Z_{k}%
)}-x^{-1/\gamma_{1}}\right\}  ,\qquad1\leq x\leq T_{n,a}.
\]
Then, for any $\varepsilon>0$,
\[
\sup_{1\leq x\leq T_{n,a}}x^{\varepsilon}\left\vert D_{n,a}^{(\beta
)}(x)-x^{-\beta/\gamma}J_{n}(x)-x^{-\beta/\gamma}x^{-1/\gamma_{1}}%
\frac{x^{\tau_{1}/\gamma_{1}}-1}{\tau_{1}\gamma_{1}}\sqrt{k}A_{1}%
(h)\right\vert =o_{\mathbf{P}}(1),
\]
as $n\rightarrow\infty$, where $J_{n}(x)$ is the Gaussian process defined in
the text.\smallskip

\noindent This approximation holds uniformly over the truncated domain $1 \le
x \le T_{n,a}$ and does not require the restriction $p>1/2$, in contrast with
the classical untruncated setting.
\end{theorem}

\noindent This result establishes that the weighted and truncated process
admits a Gaussian approximation uniformly over the truncated region, even
under strong censoring. In particular, the combination of weighting and
truncation restores a regular Gaussian-type stochastic structure that is lost
in the classical setting when $p\leq1/2$.

\begin{theorem}
[Consistency]\label{theorem2}

Assume that the survival functions $\overline{F}$ and $\overline{G}$ satisfy
the first-order regular variation conditions \eqref{RVF} and \eqref{RVG},
respectively. Let $k=k_{n}$ be an intermediate sequence such that
\[
k\to\infty\qquad\text{and} \qquad k/n\to0, \quad\text{as } n\to\infty.
\]
Let $(a_{n})$ be a positive sequence satisfying \eqref{assump_an}, and define
\[
m_{n} := \lfloor n a_{n} \rfloor.
\]
Then, for every $\beta>1$,
\[
\widehat{\gamma}_{1,k}^{(\mathrm{NA,tr})}(\beta) \overset{\mathbf{P}%
}{\longrightarrow} \gamma_{1}, \qquad\text{as } n\to\infty.
\]

\end{theorem}

\noindent This result shows that the proposed estimator is consistent for the
tail index $\gamma_{1}$ over the entire censoring range $0<p<1$, including
both weak and strong censoring regimes.

\begin{theorem}
[Asymptotic normality]\label{theorem3}

In addition to the assumptions of Theorem~\ref{theorem2}, assume that the
second-order condition \eqref{second-orderF} holds and let $h:=U_{H}(n/k)$.
Assume that
\[
\sqrt{k}A_{1}(h)=O(1),
\]
and that
\begin{equation}
\sqrt{k}\left(  \frac{na_{n}}{k}\right)  ^{\beta}\log\!\left(  \frac{k}%
{na_{n}}\right)  \rightarrow0. \label{supplementary_assumption}%
\end{equation}
Then, for any fixed $\beta>1$,
\[
\sqrt{k}\left(  \widehat{\gamma}_{1,k}^{(\mathrm{NA,tr})}(\beta)-\gamma
_{1}\right)  =\gamma_{1}N_{n}(\beta)+\frac{\beta}{\beta-p\tau_{1}}\sqrt
{k}A_{1}(h)+o_{\mathbf{P}}(1),
\]
as $n\rightarrow\infty$, where
\[
N_{n}(\beta)=\frac{\beta(p+\beta-1)}{p}\int_{0}^{1}s^{\beta-2}\mathbf{W}%
_{n,1}(s)\,ds-\frac{\beta}{p}\mathbf{W}_{n,1}(1)+\beta\int_{0}^{1}s^{\beta
-2}\mathbf{W}_{n,2}(s)\,ds.
\]
If, moreover, $\sqrt{k}A_{1}(h)\rightarrow\lambda<\infty$, then
\[
\sqrt{k}\left(  \widehat{\gamma}_{1,k}^{(\mathrm{NA,tr})}(\beta)-\gamma
_{1}\right)  \overset{\mathcal{D}}{\longrightarrow}\mathcal{N}(\mu_{\beta
},\sigma_{\beta}^{2}),
\]
where
\[
\mu_{\beta}=\frac{\lambda\beta}{\beta-p\tau_{1}},\qquad\sigma_{\beta}%
^{2}=\frac{1}{p}\frac{\beta^{2}}{2\beta-1}\gamma_{1}^{2}.
\]

\end{theorem}

\noindent The supplementary condition \eqref{supplementary_assumption} ensures
that the boundary term induced by truncation is asymptotically negligible at
the $\sqrt{k}$ scale. Consequently, the truncation has no impact on the
limiting distribution of the estimator, and the asymptotic normality is
entirely driven by the weighted tail empirical process.

\begin{remark}
A natural and practical choice of the truncation index is
\[
m_{n}=\max\left(  3,\lfloor\log\log k\rfloor\right)  .
\]
We show that this choice satisfies all the conditions required in
Theorems~\ref{theorem1}--\ref{theorem3}. First, since $k\rightarrow\infty$, we
have
\[
m_{n}=\lfloor\log\log k\rfloor\longrightarrow\infty.
\]
Next, since $\log\log k=o(k)$, it follows that
\[
\frac{m_{n}}{k}\longrightarrow0.
\]
Recall that $a_{n}=m_{n}/n$. Then
\[
a_{n}=\frac{\log\log k}{n}\longrightarrow0,
\]
since $n\rightarrow\infty$. Moreover,
\[
na_{n}=m_{n}=\log\log k\longrightarrow\infty,
\]
which verifies the condition $na_{n}\rightarrow\infty$. We now check that
$a_{n}=o(k/n)$. Indeed,
\[
\frac{a_{n}}{k/n}=\frac{m_{n}}{k}=\frac{\log\log k}{k}\longrightarrow0.
\]
Finally, we verify the strengthened condition required for the asymptotic
normality:
\[
\sqrt{k}\left(  \frac{na_{n}}{k}\right)  ^{\beta}\log\!\left(  \frac{k}%
{na_{n}}\right)  =\sqrt{k}\left(  \frac{m_{n}}{k}\right)  ^{\beta}%
\log\!\left(  \frac{k}{m_{n}}\right)  .
\]
Since $m_{n}=\log\log k$, we obtain
\[
\left(  \frac{m_{n}}{k}\right)  ^{\beta}=\left(  \frac{\log\log k}{k}\right)
^{\beta},\qquad\log\!\left(  \frac{k}{m_{n}}\right)  =\log k-\log\log\log
k\sim\log k.
\]
Hence
\[
\sqrt{k}\left(  \frac{m_{n}}{k}\right)  ^{\beta}\log\!\left(  \frac{k}{m_{n}%
}\right)  =k^{1/2-\beta}(\log\log k)^{\beta}\log k.
\]
Since $\beta>1$, we have $1/2-\beta<0$, and therefore
\[
k^{1/2-\beta}(\log\log k)^{\beta}\log k\longrightarrow0.
\]
This shows that all the required conditions are satisfied. In particular, the
choice $m_{n}=\lfloor\log\log k\rfloor$ provides a theoretically valid and
practically negligible truncation.
\end{remark}

\begin{remark}
If $p>1/2$, then $\sigma_{\beta}^{2}<\sigma_{p}^{2}$. This shows that, in the
weak censoring regime, the asymptotic variance of $\widehat{\gamma}
_{1,k}^{(\mathrm{NA,tr})}(\beta)$ (resp. $\widehat{\gamma}_{1,k}
^{(\mathrm{KM})}(\beta)$) is smaller than that of $\widehat{\gamma}
_{1,k}^{(\mathrm{MNS})}$ (resp. $\widehat{\gamma}_{1,k}^{(\mathrm{W})}$).
Hence, the proposed estimator not only remains valid under strong censoring,
but may also improve efficiency in the weak censoring regime. Indeed,
\[
\sigma_{\beta}^{2}-\sigma_{p}^{2}=\frac{(\beta-p)(2p\beta-p-\beta
)}{p(2p-1)(2\beta-1)}\gamma_{1}^{2}<0,
\]
provided that $1<\beta<p/(2p-1)$, which holds since $p/(2p-1)>1$. This
inequality provides a useful guideline for selecting $\beta$ in practice.
\end{remark}

\begin{remark}
According to \cite{BWW2019}, the asymptotic variance of $\widehat{\gamma
}_{1,k}^{(\mathrm{BW})}(\beta)$ equals $\gamma_{1}^{2}p(1+\beta\gamma_{1}%
)^{2}/(2p(1+\beta\gamma_{1})-1)$, whereas that of $\widehat{\gamma}%
_{1,k}^{(\mathrm{W})}$ is $\gamma_{1}^{2}p/(2p-1)$ for $p>1/2$. However, the
comparison between these two variances is not straightforward, as it depends
on the joint effect of $p$, $\gamma_{1}$, and $\beta$. A simple algebraic
manipulation shows that the difference between the two variances is
\[
\frac{p\beta\gamma_{1}}{2p-1}\frac{(2p-1)\beta\gamma_{1}-2(1-p)}%
{2p(1+\beta\gamma_{1})-1}.
\]
The sign of this expression is not determined a priori, which makes the
selection of $\beta$ more delicate in this framework.
\end{remark}

\begin{remark}
As previously mentioned, Theorem~\ref{theorem3} simultaneously generalizes the
asymptotic normality of the estimators $\widehat{\gamma}_{1,k}^{(\mathrm{MNS}
)}$ and $\widehat{\gamma}_{1,k}^{(\mathrm{W})}$ under the condition $p>1/2$.
More importantly, it extends these results to the full censoring range
$0<p<1$, which is not covered by existing approaches. Furthermore, it relaxes
the assumptions required to establish the asymptotic normality of
$\widehat{\gamma}_{1,k}^{(\mathrm{W})}$ in \cite{BWW2019} by extending its
validity beyond the classical Hall-type models. In addition, the theorem
provides a convenient functional representation of the limiting distribution
in terms of Wiener processes.
\end{remark}

\begin{remark}
(Asymptotic choice of the tuning parameter $\beta$) A natural way to discuss
the choice of $\beta$ is through the first-order asymptotic mean squared error
(AMSE). From Theorem~\ref{theorem3}, we have
\[
\sqrt{k}\left(  \widehat{\gamma}_{1,k}^{(\mathrm{NA,tr})}(\beta)-\gamma
_{1}\right)  =\gamma_{1}N_{n}(\beta)+\frac{\beta}{\beta-p\tau_{1}}\sqrt
{k}A_{1}(h)+o_{\mathbb{P}}(1),
\]
which yields
\[
\mathrm{AMSE}(\beta)\approx\frac{\gamma_{1}^{2}}{pk}\frac{\beta^{2}}{2\beta
-1}+\left(  \frac{\beta}{\beta-p\tau_{1}}\right)  ^{2}A_{1}(h)^{2}.
\]
For $\beta>1$, both the variance component
\[
\frac{\beta^{2}}{2\beta-1}%
\]
and the bias component
\[
\frac{\beta}{\beta-p\tau_{1}}%
\]
are increasing functions of $\beta$. Consequently, the AMSE does not admit an
interior minimizer at the first-order level.
\end{remark}

\noindent Therefore, the optimal asymptotic strategy is to select $\beta$ as
close as possible to $1$, while maintaining $\beta>1$ to preserve the validity
of the Gaussian approximation. In practice, this suggests choosing $\beta$
slightly larger than $1$, for instance $\beta=1.01$ or $\beta=1.001$.

\section{\textbf{Simulation study}\label{sec3}}

\noindent This section evaluates the finite-sample performance of the proposed
estimator $\widehat{\gamma}_{1,k}^{(\mathrm{NA,tr})}(\beta)$ and compares it
with the classical Nelson--Aalen estimator $\widehat{\gamma}_{1,k}%
^{(\mathrm{MNS})}$ and the adapted Hill estimator $\widehat{\gamma}%
_{1,k}^{(\mathrm{EFG})}$, with particular emphasis on moderate and strong
censoring regimes. The objective is to assess both the stability and accuracy
of the proposed method in settings where classical approaches are known to
deteriorate.\smallskip

\noindent To this end, we consider three standard heavy-tailed models, widely
used in extreme value analysis:

\begin{itemize}
\item the Burr distribution with parameters $(\zeta,\eta)$, defined by
\[
F(x)=1-\left(  1+x^{1/\eta}\right)  ^{-\eta/\zeta},\qquad x>0,
\]


\item the Fr\'{e}chet distribution with parameter $\zeta$, given by
\[
F(x)=\exp\!\left(  -x^{-1/\zeta}\right)  ,\qquad x>0,
\]


\item the log-gamma distribution, defined through $\log X\sim\Gamma
(\alpha,\varsigma)$, with distribution function
\[
F(x)=\frac{1}{\varsigma^{\alpha}\Gamma(\alpha)}\int_{0}^{\log x}u^{\alpha
-1}e^{-u/\varsigma}\,du,\qquad x>1.
\]

\end{itemize}

\noindent These models allow us to cover a broad range of tail behaviors,
including both regularly varying and sub-exponential-type distributions,
thereby providing a comprehensive assessment of the estimators under different
tail regimes.\smallskip

\noindent For each model, we consider the following censoring schemes:

\begin{itemize}
\item[$\bullet$] Burr$(\gamma_{1},\eta)$ censored by Burr$(\gamma_{2},\eta)$,

\item[$\bullet$] Fr\'{e}chet$(\gamma_{1})$ censored by Fr\'{e}chet$(\gamma
_{2})$,

\item[$\bullet$] Log-gamma$(\alpha_{1},\varsigma_{1})$ censored by
Log-gamma$(\alpha_{2},\varsigma_{2})$,
\end{itemize}

\noindent where $\alpha_{1}=1/\gamma_{1}$ and $\alpha_{2}=1/\gamma_{2}$. This
construction ensures that both the lifetime and censoring distributions belong
to the same parametric family, allowing for a controlled variation of the
censoring intensity.\smallskip

\noindent The model parameters are fixed as follows: $\eta=0.25$ and
$\varsigma_{1}=\varsigma_{2}=2$. We consider two values for the tail index,
$\gamma_{1}=0.4$ and $\gamma_{1}=0.7$, corresponding respectively to
moderately and strongly heavy-tailed distributions. These choices are standard
in the literature and allow us to evaluate the performance of the estimators
across different levels of tail heaviness.\smallskip

\noindent Three censoring levels are investigated: $p=0.70,\;0.50,\;\text{and
}0.30$, representing weak, moderate, and strong censoring regimes,
respectively. These values are chosen to explicitly assess the performance of
the estimators across the full censoring spectrum, including the challenging
regime $p\leq1/2$, where classical methods typically fail. For each pair
$(\gamma_{1},p)$, the corresponding value of $\gamma_{2}$ is obtained from
$p=\frac{\gamma_{2}}{\gamma_{1}+\gamma_{2}}$, ensuring consistency between the
tail behavior of the signal and the censoring distribution.\smallskip

\noindent The effect of the tuning parameter is studied for $\beta
=1.01,\;1.5,\;\text{and }2$. The value $\beta=1.01$ corresponds to the
theoretically recommended regime, where $\beta$ is chosen close to $1$. This
choice is motivated by the asymptotic theory, which suggests that values of
$\beta$ close to $1$ provide a favorable balance between bias and variance
while ensuring stability under strong censoring. The additional values
$\beta=1.5$ and $\beta=2$ are included to assess the impact of stronger
weighting on the bias--variance trade-off.\smallskip

\noindent For each configuration, $2000$ independent samples of size $n=1000$
are generated. The sample size $n=1000$ is standard in extreme value
simulations and ensures that a sufficiently large number of upper order
statistics is available for reliable estimation. This choice provides a
reasonable compromise between computational cost and statistical
accuracy.\smallskip

\noindent The results are summarized in Figures~\ref{fig1}--\ref{fig9}
(Appendix~B), which report the empirical bias and mean squared error (MSE) of
the estimators as functions of the number of upper order statistics $k$. This
representation allows us to analyze both the stability of the estimators with
respect to the threshold choice and their overall efficiency across different
censoring regimes.

\subsection{Discussion of the simulation results}

\noindent The simulation results reported in Figures~\ref{fig1}--\ref{fig9}
provide a systematic comparison between the proposed estimator
$\widehat{\gamma}_{1,k}^{(\mathrm{NA,tr})}(\beta)$ and the competing
estimators $\widehat{\gamma}_{1,k}^{(\mathrm{MNS})}$ and $\widehat{\gamma
}_{1,k}^{(\mathrm{EFG})}$ across different censoring levels and heavy-tailed
models. The analysis focuses on both accuracy (bias) and efficiency (MSE), as
well as on the stability with respect to the threshold parameter
$k$.\smallskip

\noindent A first key observation is that the performance of all estimators
strongly depends on the censoring intensity. Under strong censoring
($p=0.30$), the proposed estimator clearly outperforms the Nelson--Aalen
estimator $\widehat{\gamma}_{1,k}^{(\mathrm{MNS})}$ in terms of both bias and
mean squared error (MSE), especially for moderate values of $k$. This behavior
reflects the instability of classical methods in this regime, where the
effective amount of uncensored tail information is limited. In particular, the
improvement is most visible in the region where classical estimators exhibit
high variability. This improvement is particularly pronounced for $\beta$
close to $1$, confirming the theoretical recommendation. In contrast, larger
values of $\beta$ (e.g., $\beta=2$) tend to increase the bias as $k$ grows,
reflecting the bias--variance trade-off induced by the weighting
scheme.\smallskip

\noindent Under moderate censoring ($p=0.50$), the performance gap between
estimators narrows, as more uncensored observations become available in the
tail. Nevertheless, the proposed estimator with $\beta$ close to $1$ maintains
a stable behavior, with bias values close to zero and consistently smaller MSE
over a wide range of $k$. This indicates that the method adapts smoothly to
intermediate censoring regimes.\smallskip

\noindent In the weak censoring regime ($p=0.70$), all estimators exhibit
similar performance, which is expected since censoring has a limited impact on
the tail. In this setting, the proposed estimator remains competitive, showing
no noticeable loss of efficiency compared to classical methods. This confirms
that the proposed modification does not degrade performance in favorable
scenarios.\smallskip

\noindent These conclusions are remarkably consistent across the three
considered models (Burr, Fr\'{e}chet, and log-gamma), indicating that the
proposed methodology is robust with respect to the underlying distribution. In
particular, no significant model-dependent degradation is observed. This
robustness highlights the general applicability of the approach.\smallskip

\noindent Another important aspect concerns the role of the tuning parameter
$\beta$. The simulations clearly show that values of $\beta$ slightly larger
than $1$ provide the best compromise between bias and variance. This
observation is fully consistent with the asymptotic analysis, which indicates
that the optimal choice of $\beta$ lies close to $1$. Conversely, larger
values of $\beta$ tend to overweight extreme observations, resulting in
increased bias for large values of $k$. This confirms the practical relevance
of the theoretical guidance for selecting $\beta$.\smallskip

\noindent As expected in extreme value estimation, all estimators exhibit a
U-shaped MSE curve as a function of $k$, reflecting the classical
bias--variance trade-off. However, the proposed estimator with $\beta$ close
to $1$ provides a wider stability region, with relatively flat MSE curves over
a broad range of $k$. This feature is particularly valuable in practice, as it
reduces sensitivity to the choice of the threshold parameter. In particular,
it alleviates the need for fine tuning of $k$.\smallskip

\noindent Finally, the most significant improvement brought by the proposed
estimator is observed under strong censoring. In this regime, classical
estimators may suffer from instability and increased variability, whereas the
proposed weighted and truncated estimator remains stable and accurate. This
clearly illustrates the advantage of the proposed regularization mechanism.
This empirical behavior is consistent with the theoretical results established
in Section~\ref{sec2}, which guarantee asymptotic validity over the entire
censoring range $0<p<1$.\smallskip

\noindent Overall, the simulation study confirms that the combination of
weighting and truncation provides a substantial improvement in finite-sample
performance, particularly in challenging censoring scenarios, while preserving
efficiency in favorable settings, including the weak censoring regime. These
findings strongly support the practical relevance of the proposed methodology.

\section{\textbf{Real data application\label{sec4}}}

\subsection{Insurance loss data}

\noindent The insurance loss dataset, collected by the US Insurance Services
Office, Inc., is publicly available through the copula package in the R
statistical software. This dataset has been extensively studied in the
literature; see, for instance, \cite{FV98}, \cite{Klugman99}, and
\cite{DPV06}. It consists of $1500$ observations, among which $34$ are
right-censored, indicating that the corresponding losses exceed policy limits
that vary across contracts.\smallskip

\noindent Such datasets are typical in actuarial science and risk management,
where large insurance claims are often subject to policy limits or reporting
thresholds. As a result, the exact magnitude of extreme losses is not fully
observed, but only known to exceed a given bound. This naturally leads to
right-censored heavy-tailed data, for which appropriate statistical methods
are required to ensure reliable inference in the tail. In particular, this
type of censoring mechanism directly motivates the use of estimators that
remain stable under incomplete tail information.\smallskip

\noindent Since our analysis is conducted within the framework of Pareto-type
tail models, it is essential to assess whether the data are compatible with
this class of distributions. To this end, we consider the log--log survival
plot based on the Nelson--Aalen estimator $F^{(\mathrm{NA})}$, which properly
accounts for right-censored observations.\smallskip

\noindent Under a Pareto-type tail behavior, the survival function satisfies
\[
\mathbb{P}(X>x)\sim Cx^{-1/\gamma_{1}},\qquad x\rightarrow\infty,
\]
which implies that the plot of $\log\overline{F}^{(\mathrm{NA})}(x)$ versus
$\log x$ should be approximately linear in the upper tail. This diagnostic is
standard in extreme value analysis and provides a graphical validation of the
regular variation assumption.\smallskip

\noindent The resulting plot (Figure~\ref{loss}) exhibits a clear linear trend
for the largest observations, with a negative slope, providing empirical
support for a Pareto-type tail behavior. Although deviations from linearity
are observed for smaller values, the upper tail remains well aligned with the
theoretical model. This pattern is typical of heavy-tailed data, where the
Pareto approximation is expected to hold only beyond a sufficiently high
threshold. This empirical evidence supports the use of extreme value methods
for estimating the tail index and justifies the application of the proposed
estimator in this setting.\smallskip

\noindent The log--log survival plot (Figure~\ref{loss}) provides empirical
support for the use of extreme value methods in this context. In particular,
the approximate linearity observed in the upper tail, together with the
negative slope, is consistent with a Pareto-type tail behavior. This suggests
that the regular variation assumption is reasonable for sufficiently large
observations. Consequently, the use of tail index estimators adapted to
heavy-tailed distributions under right censoring is well justified in this
setting.
%TCIMACRO{\FRAME{ftbpFU}{4.1831in}{2.6394in}{0pt}{\Qcb{Log--log survival plot
%based on the Nelson--Aalen estimator for the insurance loss data. The
%approximate linearity observed in the upper tail, along with the negative
%slope, is characteristic of a Pareto-type heavy-tailed behavior.}}{\Qlb{loss}%
%}{loss-1.eps}{\special{ language "Scientific Word";  type "GRAPHIC";
%maintain-aspect-ratio TRUE;  display "USEDEF";  valid_file "F";
%width 4.1831in;  height 2.6394in;  depth 0pt;  original-width 9.3054in;
%original-height 5.8496in;  cropleft "0";  croptop "1";  cropright "1";
%cropbottom "0";  filename 'loss-1.eps';file-properties "XNPEU";}} }%
%BeginExpansion
\begin{figure}[ptb]%
\centering
\includegraphics[
height=2.6394in,
width=4.1831in
]%
{loss-1.eps}%
\caption{Log--log survival plot based on the Nelson--Aalen estimator for the
insurance loss data. The approximate linearity observed in the upper tail,
along with the negative slope, is characteristic of a Pareto-type heavy-tailed
behavior.}%
\label{loss}%
\end{figure}
%EndExpansion


\noindent Using the adaptive selection procedure of Reiss and Thomas
\citep{ReTo97}, the optimal sample fraction for estimating the censoring
proportion is $k_{\mathrm{opt}}=51$, yielding $\widehat{p}_{k}=0.76$. This
relatively large value confirms that the dataset falls within the weak
censoring regime, where classical estimators are expected to perform
reasonably well.\smallskip

\noindent The adapted Hill estimator $\widehat{\gamma}_{1,k}^{(\mathrm{EFG})}$
achieves its optimal value at $k=73$, leading to the estimate $\widehat{\gamma
}_{1,k}^{(\mathrm{EFG})}=0.77$. For the Nelson--Aalen tail index estimator
$\widehat{\gamma}_{1,k}^{(\mathrm{MNS})}$, the optimal choice is $k=30$,
resulting in $\widehat{\gamma}_{1,k}^{(\mathrm{MNS})}=0.45$. \smallskip

\noindent The same selection procedure is applied to the proposed estimator.
Following the theoretical recommendations of Section~\ref{sec3}, we set
$\beta=1.01$. The corresponding optimal number of upper order statistics is
$k=30$, yielding the estimate $\widehat{\gamma}_{1,k}^{(\mathrm{NA,tr})}%
(\beta)=0.51$.\smallskip

\noindent The adaptive selection of $k$ is based on the criterion
\begin{equation}
k_{\mathrm{opt}}:=\arg\min_{1<k<n}\frac{1}{k}\sum_{i=1}^{k}i^{\nu}\left\vert
\widehat{\xi}_{i}-\mathrm{median}\left(  \widehat{\xi}_{1},\ldots
,\widehat{\xi}_{k}\right)  \right\vert ,\qquad0\leq\nu\leq\frac{1}{2},
\label{kopt}%
\end{equation}
where $\widehat{\xi}_{i}$ denotes a generic tail index estimator computed from
the $i$ largest observations. Following \cite{NF2004}, we set $\nu=0.3$, which
provides a good compromise between bias and variability.\smallskip

\noindent A notable feature of the results is that the proposed estimator
yields an estimate that lies between those obtained from the adapted Hill
estimator and the classical Nelson--Aalen estimator. This intermediate
behavior is consistent with its theoretical construction, which balances bias
and variance through the tuning parameter $\beta$.\smallskip

\noindent More importantly, the proposed estimator provides a stable estimate
while relying on a relatively moderate number of upper order statistics
($k=30$), indicating reduced sensitivity to the threshold selection. This
reduced sensitivity is particularly desirable in practice, where the choice of
$k$ is often a major source of uncertainty. This observation is in line with
the simulation results of Section~\ref{sec3}, where the estimator exhibits a
wide stability region with respect to $k$.\smallskip

\noindent Overall, the empirical analysis confirms that the proposed weighted
and truncated estimator provides a reliable and stable estimation of the tail
index in real data applications, while remaining consistent with both
theoretical predictions and simulation evidence.

\subsection{The Australian AIDS survival dataset}

\noindent The Australian AIDS survival dataset, compiled by Dr. P. J. Solomon
and the Australian National Centre for HIV Epidemiology and Clinical Research,
consists of $2843$ patients diagnosed prior to July 1, 1991. In this study, we
focus on a subset of $2754$ male patients. This dataset has been extensively
analyzed in the literature; see, for instance, \cite{RS-94}, \cite{VR-02}
(pp.~379--385), \cite{EnFG08}, and \cite{Ndao16}.\smallskip

\noindent The observations correspond to survival times measured from AIDS
diagnosis until death or the end of the study, leading naturally to
right-censored data. Such datasets are typical in survival analysis and
provide a challenging framework for extreme value methods due to the potential
scarcity of uncensored observations in the tail. In particular, this setting
is representative of moderate to strong censoring regimes, where the effective
amount of tail information may be severely limited.\smallskip

\noindent As in the previous application, we assess the compatibility of the
data with a Pareto-type tail by examining the log--log survival plot based on
the Nelson--Aalen estimator. The plot, displayed in Figure~\ref{aids-1},
exhibits an approximately linear behavior in the upper tail together with a
negative slope, which is consistent with a heavy-tailed Pareto-type
distribution. This suggests that the regular variation assumption is
reasonable for sufficiently large survival times. Although deviations from
linearity are visible over intermediate ranges, the tail behavior remains
sufficiently regular to justify the use of extreme value techniques.
%TCIMACRO{\FRAME{ftbpFU}{4.1831in}{2.6394in}{0pt}{\Qcb{Log--log survival plot
%based on the Nelson--Alen estimator for the Australian AIDS survival data. The
%approximately linear behavior observed in the upper tail, together with the
%negative slope, indicates compatibility with a Pareto-type heavy-tailed
%model.}}{\Qlb{aids-1}}{aids-1.eps}{\special{ language "Scientific Word";
%type "GRAPHIC";  maintain-aspect-ratio TRUE;  display "USEDEF";
%valid_file "F";  width 4.1831in;  height 2.6394in;  depth 0pt;
%original-width 9.3054in;  original-height 5.8496in;  cropleft "0";
%croptop "1";  cropright "1";  cropbottom "0";
%filename 'aids-1.eps';file-properties "XNPEU";}} }%
%BeginExpansion
\begin{figure}[ptb]%
\centering
\includegraphics[
height=2.6394in,
width=4.1831in
]%
{aids-1.eps}%
\caption{Log--log survival plot based on the Nelson--Alen estimator for the
Australian AIDS survival data. The approximately linear behavior observed in
the upper tail, together with the negative slope, indicates compatibility with
a Pareto-type heavy-tailed model.}%
\label{aids-1}%
\end{figure}
%EndExpansion


\noindent Using the adaptive selection procedure of Reiss and Thomas
\citep{ReTo97}, the optimal sample fraction for estimating the censoring
proportion is $k_{\mathrm{opt}}=162$, yielding $\widehat{p}_{k}=0.30$. This
clearly indicates that the dataset falls within the strong censoring regime,
where classical asymptotic results are no longer guaranteed. \smallskip

\noindent For the adapted Hill estimator, the optimal threshold is $k=55$,
leading to the estimate $\widehat{\gamma}_{1,k}^{(\mathrm{EFG})}=0.72$. The
same threshold is selected for the classical Nelson--Aalen estimator
$\widehat{\gamma}_{1,k}^{(\mathrm{MNS})}$, but the resulting estimate is
significantly lower, $\widehat{\gamma}_{1,k}^{(\mathrm{MNS})}=0.15$.
\smallskip

\noindent This pronounced discrepancy highlights the severe instability of
classical estimators under strong censoring, where the effective number of
uncensored extreme observations becomes too small to ensure reliable
inference. In particular, the Nelson--Aalen estimator exhibits a substantial
downward bias in this setting.\smallskip

\noindent In contrast, applying the proposed estimator with $\beta=1.01$
yields a substantially larger optimal threshold, $k=275$, and the estimate
$\widehat{\gamma}_{1,k}^{(\mathrm{NA,tr})}(\beta)=0.64$.\smallskip

\noindent Two important features emerge from this result. First, the estimator
remains stable despite the strong censoring level, producing a value that is
coherent with the adapted Hill estimator while avoiding the severe
underestimation observed for the classical Nelson--Aalen estimator. This
confirms that the proposed method effectively corrects the bias induced by
censoring. Second, the larger selected value of $k$ indicates that the
proposed method is able to exploit a wider portion of the upper tail,
reflecting its improved robustness with respect to censoring. This ability to
use more extreme observations is a key practical advantage.\smallskip

\noindent These observations are fully consistent with the theoretical results
established in Section~\ref{sec2}, which guarantee the validity of the
estimator over the entire censoring range $0<p<1$. In particular, the
combination of weighting and truncation effectively stabilizes the tail
empirical process, even when the proportion of uncensored observations is
small.\smallskip

\noindent Overall, this example clearly illustrates the main advantage of the
proposed estimator: it provides reliable and stable tail index estimation
precisely in the regime where classical methods break down.

\section{\textbf{Proofs}\label{sec5}}

\subsection{Proof of Theorem \ref{theorem1}}

\noindent The proof follows the Gaussian approximation scheme developed in
\cite{MNS2025} for the Nelson--Aalen tail process. We recall only the
modifications required by the truncation and the weighting.\smallskip

\noindent Let
\[
D_{n}(x):=\sqrt{k}\left\{  \frac{\overline{F}_{n}^{(\mathrm{NA})}(xZ_{n-k:n}%
)}{\overline{F}_{n}^{(\mathrm{NA})}(Z_{n-k:n})}-x^{-1/\gamma_{1}}\right\}
,\qquad x\geq1.
\]
Under the first- and second-order regular variation assumptions, the classical
decomposition of the Nelson--Aalen tail process gives
\[
D_{n}(x)=J_{n}(x)+x^{-1/\gamma_{1}}\frac{x^{\tau_{1}/\gamma_{1}}-1}{\tau
_{1}\gamma_{1}}\sqrt{k}A_{1}(h)+R_{n}(x),
\]
where $J_{n}(x)$ is the Gaussian term defined in the statement of the theorem
and $R_{n}(x)$ is the corresponding remainder.\smallskip

\noindent In the untruncated case, the control of $R_{n}(x)$ involves the
lower boundary of the transformed empirical process and produces a boundary
contribution of order $k^{1/2-p}$, which is negligible only when $p>1/2$. In
the present setting, the analysis is restricted to
\[
1\leq x\leq T_{n,a}:=\frac{u_{n}}{Z_{n-k:n}},
\]
where $u_{n}=Q^{(1)}(a_{n})$ and
\[
a_{n}\rightarrow0,\qquad na_{n}\rightarrow\infty,\qquad a_{n}=o(k/n).
\]
Hence the transformed domain is bounded away from zero by a quantity of order
$a_{n}$. Consequently, the lower-boundary contribution appearing in the
untruncated proof is removed. The same decomposition therefore yields, for
some $1/4<\eta<1/2$ and any sufficiently small $\varepsilon_{0}>0$,
\[
R_{n}(x)=o_{\mathbf{P}}\!\left(  x^{(2\eta-p)/\gamma+\varepsilon_{0}}\right)
,\qquad\text{uniformly for }1\leq x\leq T_{n,a}.
\]
Now define the weighted and truncated process by
\[
D_{n,a}^{(\beta)}(x)=x^{-\beta/\gamma}D_{n}(x),\qquad1\leq x\leq T_{n,a}.
\]
Multiplying the preceding expansion by $x^{-\beta/\gamma}$ gives
\[
D_{n,a}^{(\beta)}(x)=x^{-\beta/\gamma}J_{n}(x)+x^{-\beta/\gamma}%
x^{-1/\gamma_{1}}\frac{x^{\tau_{1}/\gamma_{1}}-1}{\tau_{1}\gamma_{1}}\sqrt
{k}A_{1}(h)+x^{-\beta/\gamma}R_{n}(x).
\]
It remains to control the last term. Since
\[
R_{n}(x)=o_{\mathbf{P}}\!\left(  x^{(2\eta-p)/\gamma+\varepsilon_{0}}\right)
\]
uniformly on the truncated domain, we have
\[
x^{-\beta/\gamma}R_{n}(x)=o_{\mathbf{P}}\!\left(  x^{(2\eta-p-\beta
)/\gamma+\varepsilon_{0}}\right)  .
\]
Choose $\eta$ and $\varepsilon_{0}$ such that
\[
\frac{2\eta-p-\beta}{\gamma}+\varepsilon_{0}<0.
\]
This is possible because $\beta>1$ and $2\eta<1$. Hence, for any sufficiently
small $\varepsilon>0$,
\[
\sup_{1\leq x\leq T_{n,a}}x^{\varepsilon}\left\vert x^{-\beta/\gamma}%
R_{n}(x)\right\vert =o_{\mathbf{P}}(1).
\]
Therefore,
\[
\sup_{1\leq x\leq T_{n,a}}x^{\varepsilon}\left\vert D_{n,a}^{(\beta
)}(x)-x^{-\beta/\gamma}J_{n}(x)-x^{-\beta/\gamma}x^{-1/\gamma_{1}}%
\frac{x^{\tau_{1}/\gamma_{1}}-1}{\tau_{1}\gamma_{1}}\sqrt{k}A_{1}%
(h)\right\vert =o_{\mathbf{P}}(1).
\]
This completes the proof.

\subsection{Proof of Theorem \ref{theorem2}}

\noindent We prove the consistency of the weighted and truncated Nelson--Aalen
estimator $\widehat{\gamma}_{1,k}^{(\mathrm{NA,tr})}(\beta)$. Set
\[
Z_{k}:=Z_{n-k:n},\qquad T_{n,a}:=\frac{u_{n}}{Z_{k}},\qquad\alpha_{n}%
:=\frac{\beta}{\widehat{p}_{k}},
\]
and define, for $1\leq x\leq T_{n,a}$,
\[
Q_{n}(x):=\frac{\overline{F}_{n}^{(\mathrm{NA})}(xZ_{k})}{\overline{F}%
_{n}^{(\mathrm{NA})}(Z_{k})}.
\]
By Proposition~\ref{prop:app-functional},
\begin{equation}
\widehat{\gamma}_{1,k}^{(\mathrm{NA,tr})}(\beta)=\alpha_{n}\int_{1}^{T_{n,a}%
}x^{-1}\left(  Q_{n}(x)\right)  ^{\alpha_{n}}\,dx-\alpha_{n}(\log
T_{n,a})\left(  Q_{n}(T_{n,a})\right)  ^{\alpha_{n}}.
\label{eq:consistency-representation}%
\end{equation}
Since $\widehat{p}_{k}\overset{\mathbf{P}}{\longrightarrow}p$, we have
\[
\alpha_{n}=\frac{\beta}{\widehat{p}_{k}}\overset{\mathbf{P}}{\longrightarrow
}\frac{\beta}{p}.
\]
For each fixed $x\geq1$, Lemma~\ref{lemma:power-convergence-alpha} gives
\[
\left(  Q_{n}(x)\right)  ^{\alpha_{n}}\overset{\mathbf{P}}{\longrightarrow
}x^{-\beta/\gamma}.
\]
Moreover, by Lemma~\ref{lem:domination-Qn}, for any sufficiently small
$\varepsilon_{0}>0$, there exists a constant $C>0$ such that, with probability
tending to one,
\[
x^{-1}\left(  Q_{n}(x)\right)  ^{\alpha_{n}}\leq Cx^{-1-\beta/\gamma
+\varepsilon_{0}},\qquad1\leq x\leq T_{n,a}.
\]
Choosing $\varepsilon_{0}$ sufficiently small, the function $x\mapsto
x^{-1-\beta/\gamma+\varepsilon_{0}}$ is integrable on $[1,\infty)$. Therefore,
by a standard truncation argument combined with dominated convergence on
compact intervals,
\[
\alpha_{n}\int_{1}^{T_{n,a}}x^{-1}\left(  Q_{n}(x)\right)  ^{\alpha_{n}%
}\,dx\overset{\mathbf{P}}{\longrightarrow}\frac{\beta}{p}\int_{1}^{\infty
}x^{-1-\beta/\gamma}\,dx.
\]
Since
\[
\frac{\beta}{p}\int_{1}^{\infty}x^{-1-\beta/\gamma}\,dx=\frac{\beta}{p}%
\cdot\frac{\gamma}{\beta}=\frac{\gamma}{p}=\gamma_{1},
\]
we obtain
\begin{equation}
\alpha_{n}\int_{1}^{T_{n,a}}x^{-1}\left(  Q_{n}(x)\right)  ^{\alpha_{n}%
}\,dx\overset{\mathbf{P}}{\longrightarrow}\gamma_{1}.
\label{eq:consistency-main}%
\end{equation}


It remains to control the boundary term. Since
\[
Q_{n}(T_{n,a})=\frac{\overline{F}_{n}^{(\mathrm{NA})}(u_{n})}{\overline{F}%
_{n}^{(\mathrm{NA})}(Z_{k})},
\]
Potter-type bounds and the consistency of the Nelson--Aalen tail ratio imply
\[
Q_{n}(T_{n,a})=O_{\mathbf{P}}\!\left[  \left(  \frac{na_{n}}{k}\right)
^{p}\right]  .
\]
Since $\alpha_{n}\overset{\mathbf{P}}{\longrightarrow}\beta/p$, it follows
that
\[
\left(  Q_{n}(T_{n,a})\right)  ^{\alpha_{n}}=O_{\mathbf{P}}\!\left[  \left(
\frac{na_{n}}{k}\right)  ^{\beta}\right]  .
\]
Furthermore,
\[
\log T_{n,a}=O_{\mathbf{P}}\!\left[  \log\!\left(  \frac{k}{na_{n}}\right)
\right]  .
\]
Thus,
\[
\alpha_{n}(\log T_{n,a})\left(  Q_{n}(T_{n,a})\right)  ^{\alpha_{n}%
}=O_{\mathbf{P}}\!\left[  \left(  \frac{na_{n}}{k}\right)  ^{\beta}%
\log\!\left(  \frac{k}{na_{n}}\right)  \right]  =o_{\mathbf{P}}(1),
\]
because $na_{n}/k\rightarrow0$ and the logarithmic factor is dominated by any
negative power of $na_{n}/k$. Hence,
\begin{equation}
\alpha_{n}(\log T_{n,a})\left(  Q_{n}(T_{n,a})\right)  ^{\alpha_{n}%
}=o_{\mathbf{P}}(1). \label{eq:consistency-boundary}%
\end{equation}
Combining \eqref{eq:consistency-representation}, \eqref{eq:consistency-main},
and \eqref{eq:consistency-boundary}, we conclude that
\[
\widehat{\gamma}_{1,k}^{(\mathrm{NA,tr})}(\beta)\overset{\mathbf{P}%
}{\longrightarrow}\gamma_{1}.
\]
This completes the proof.

\subsection{\textbf{Proof of Theorem} \ref{theorem3}}

\noindent By Proposition~\ref{prop:linearization}, we have
\[
\sqrt{k}\left(  \widehat{\gamma}_{1,k}^{(\mathrm{NA,tr})}(\beta)-\gamma
_{1}\right)  =\left(  \frac{\beta}{p}\right)  ^{2}\int_{1}^{T_{n,a}%
}x^{-1+1/\gamma_{1}}D_{n,a}^{(\beta)}(x)\,dx+B_{n}+o_{\mathbf{P}}(1),
\]
where
\[
B_{n}=\sqrt{k}\left[  \frac{\beta}{p}\int_{1}^{T_{n,a}}x^{-1}\left(
\frac{\overline{F}(Z_{k}x)}{\overline{F}(Z_{k})}\right)  ^{\beta/p}%
dx-\gamma_{1}\right]  .
\]
By the second-order regular variation condition \eqref{second-orderF},
\[
B_{n}=\frac{\beta}{\beta-p\tau_{1}}\sqrt{k}A_{1}(h)+o_{\mathbf{P}}(1),\qquad
h=U_{H}(n/k).
\]
It remains to analyze the stochastic integral. By Theorem~\ref{theorem1},
uniformly for $1\leq x\leq T_{n,a}$,
\[
D_{n,a}^{(\beta)}(x)=x^{-\beta/\gamma}J_{n}(x)+x^{-\beta/\gamma}%
x^{-1/\gamma_{1}}\frac{x^{\tau_{1}/\gamma_{1}}-1}{\tau_{1}\gamma_{1}}\sqrt
{k}A_{1}(h)+r_{n}(x),
\]
where
\[
\sup_{1\leq x\leq T_{n,a}}x^{\varepsilon}|r_{n}(x)|=o_{\mathbf{P}}(1).
\]
The contribution of $r_{n}$ to the integral is therefore $o_{\mathbf{P}}%
(1)$.\smallskip

\noindent Consequently,
\[
\left(  \frac{\beta}{p}\right)  ^{2}\int_{1}^{T_{n,a}}x^{-1+1/\gamma_{1}%
}D_{n,a}^{(\beta)}(x)\,dx=\left(  \frac{\beta}{p}\right)  ^{2}\int%
_{1}^{T_{n,a}}x^{-1-(\beta-p)/\gamma}J_{n}(x)\,dx+o_{\mathbf{P}}(1).
\]
By Proposition~\ref{prop:app-gaussian},
\[
\left(  \frac{\beta}{p}\right)  ^{2}\int_{1}^{T_{n,a}}x^{-1-(\beta-p)/\gamma
}J_{n}(x)\,dx=\gamma_{1}\mathcal{N}_{n}(\beta)+o_{\mathbf{P}}(1),
\]
where
\[
\mathcal{N}_{n}(\beta)=\frac{\beta(p+\beta-1)}{p}\int_{0}^{1}s^{\beta
-2}\mathbf{W}_{n,1}(s)\,ds+\beta\int_{0}^{1}s^{\beta-2}\mathbf{W}%
_{n,2}(s)\,ds-\frac{\beta}{p}\mathbf{W}_{n,1}(1).
\]
Moreover,
\[
\mathcal{N}_{n}(\beta)\overset{\mathcal{D}}{\longrightarrow}\mathcal{N}\left(
0,\frac{1}{p}\frac{\beta^{2}}{2\beta-1}\right)  .
\]
Combining the preceding estimates, we obtain
\[
\sqrt{k}\left(  \widehat{\gamma}_{1,k}^{(\mathrm{NA,tr})}(\beta)-\gamma
_{1}\right)  =\gamma_{1}\mathcal{N}_{n}(\beta)+\frac{\beta}{\beta-p\tau_{1}%
}\sqrt{k}A_{1}(h)+o_{\mathbf{P}}(1).
\]
If $\sqrt{k}A_{1}(h)\rightarrow\lambda$, then Slutsky's lemma yields
\[
\sqrt{k}\left(  \widehat{\gamma}_{1,k}^{(\mathrm{NA,tr})}(\beta)-\gamma
_{1}\right)  \overset{\mathcal{D}}{\longrightarrow}\mathcal{N}\left(
\mu_{\beta},\sigma_{\beta}^{2}\right)  ,
\]
where
\[
\mu_{\beta}=\frac{\lambda\beta}{\beta-p\tau_{1}},\qquad\sigma_{\beta}%
^{2}=\frac{1}{p}\frac{\beta^{2}}{2\beta-1}\gamma_{1}^{2}.
\]
This completes the proof.

\section{\textbf{Conclusion}\label{sec6}}

\noindent This paper introduces a new class of weighted and truncated tail
index estimators under random right censoring, based on a modified
Nelson--Aalen integral representation. The proposed estimator $\widehat{\gamma
}_{1,k}^{(\mathrm{NA,tr})}(\beta)$ is specifically designed to overcome a
fundamental limitation of existing approaches, namely their restriction to the
weak censoring regime $p>1/2$. In particular, the methodology addresses a
longstanding gap in the literature by providing a theoretically sound and
practically implementable solution under arbitrary censoring levels.
\smallskip

\noindent A central aspect of the proposed approach is the introduction of a
weighted and truncated Nelson--Aalen tail empirical process, which constitutes
the fundamental object of the analysis. The estimator is constructed as a
functional of this process, allowing its asymptotic behavior to be directly
derived from the stochastic properties of the underlying empirical structure.
This process-level perspective represents a conceptual shift from classical
approaches, which typically focus on modifying the estimator itself rather
than its underlying stochastic representation. \smallskip

\noindent The main theoretical contribution lies in establishing a uniform
Gaussian approximation for this weighted and truncated tail empirical process
over the entire censoring range $0<p<1$. This is achieved through a principled
combination of truncation and weighting: truncation removes the pathological
contribution arising from the lower part of the transformed scale, while
weighting ensures uniform control of stochastic fluctuations. As a
consequence, the proposed estimator is consistent and asymptotically normal
for all censoring levels, thereby extending classical extreme value theory
beyond the weak censoring framework. Importantly, this result shows that the
loss of asymptotic normality under strong censoring is not intrinsic, but can
be overcome through an appropriate regularization of the tail empirical
process. \smallskip

\noindent A key feature of the approach is that, under standard first- and
second-order regular variation conditions on $F$ and the usual intermediate
sequence conditions $k\rightarrow\infty$ and $k/n\rightarrow0$, the estimator
admits a linearization as a functional of the weighted and truncated tail
process. The truncation level is chosen through a sequence $m_{n}$ satisfying
$m_{n}\rightarrow\infty$ and $m_{n}/k\rightarrow0$, with typical choices
including slowly varying rates such as $m_{n}=\lfloor\log\log k\rfloor$ or
polynomial rates $m_{n}=\lfloor k^{\varepsilon}\rfloor$ with $0<\varepsilon
<1$. This representation provides a transparent connection between the
stochastic structure of the process and the asymptotic behavior of the
estimator, and plays a central role in establishing both consistency and
asymptotic normality. In particular, the use of extremely slow truncation
rates ensures that the method remains practically negligible while being
asymptotically essential. \smallskip

\noindent From a methodological perspective, the proposed estimator offers a
unified treatment of weak and strong censoring regimes. In particular,
choosing the tuning parameter $\beta$ close to $1$ (e.g., $\beta=1.01$) yields
a stable and practically effective procedure while preserving the theoretical
guarantees. In the weak censoring case, the estimator is asymptotically
equivalent to the classical Nelson--Aalen estimator, ensuring full
compatibility with existing methods. This highlights that the proposed
approach does not replace classical estimators, but rather extends them in a
natural and coherent way. \smallskip

\noindent The simulation study confirms the theoretical findings and shows
that the proposed estimator performs favorably across a wide range of models
and censoring scenarios. In particular, it exhibits improved stability and
reduced mean squared error compared to existing Nelson--Aalen-based
estimators, especially under moderate and strong censoring. These results are
further supported by real data applications, where the estimator provides
reliable tail index estimates in challenging settings. This empirical evidence
confirms that the theoretical advantages of the method translate into
practical performance gains. \smallskip

\noindent From an asymptotic efficiency viewpoint, the estimator achieves a
smaller limiting variance than the classical Nelson--Aalen estimator under
weak censoring, while remaining applicable over the full censoring range.
Although its variance may exceed that of the adapted Hill estimator, its
overall performance in terms of mean squared error is often superior,
reflecting a favorable bias--variance trade-off. This illustrates the
effectiveness of the weighting mechanism in balancing robustness and
efficiency. \smallskip

\noindent The theoretical results are established under general first- and
second-order regular variation conditions, without relying on restrictive
parametric assumptions such as Hall-type models. This ensures applicability to
a broad class of heavy-tailed distributions encountered in practice.
Consequently, the proposed framework retains a high level of generality while
achieving strong asymptotic guarantees. \smallskip

\noindent Several directions for future research naturally arise. In
particular, incorporating bias-reduction techniques within the present
framework could further improve performance. Another promising direction is
the development of data-driven procedures for the joint selection of the
tuning parameter $\beta$ and the threshold $k$. More broadly, the weighted and
truncated tail empirical process introduced in this paper provides a flexible
basis for further developments in extreme value analysis under censoring.
Extensions to dependent data or multivariate extreme value settings also
constitute promising directions for future investigation. \smallskip

\noindent In summary, the proposed approach shows that controlling the tail
empirical process itself, rather than only the estimator, provides a powerful
and flexible strategy for handling censoring in extreme value theory. The
joint use of weighting and truncation restores a unified and robust asymptotic
framework under random censoring, effectively bridging the gap between
theoretical validity and practical applicability. This work thus opens the way
for a new class of process-based methodologies in censored extreme value analysis.

\section{\textbf{Declarations}}

\begin{itemize}
\item The authors have no findings to declare.

\item The authors have no relevant financial or non-financial interests to disclose.

\item The authors have no competing interests, that are relevant to the
content of this article, to declare.

\item All authors certify that they have no affiliations with or involvement
in any organization or entity with any financial interest or non-financial
interest in the subject matter or materials discussed in this manuscript.

\item The authors have no financial or proprietary interests in any material
discussed in this article.
\end{itemize}

\renewcommand{\refname}{\textbf{REFERENCES}}

\begin{thebibliography}{99999999999999999999999999999999999999}                                                           %


\bibitem[Basu \textit{et al.}(1998)]{Basu98}Basu, A., Harris, I. R., Hjort, N.
L., \& Jones, M. C. 1998. Robust and efficient estimation by minimizing a
density power divergence. Biometrika, 85, 549--559.

\bibitem[Beirlant \textit{et al.}(2016)]{BBWG2016}Beirlant, J., Bardoutsos,
A., de Wet, T., Gijbels, I., 2016. Bias reduced tail estimation for censored
Pareto type distributions. Statist. Probab. Lett. \textbf{109}, 78-88.

\bibitem[Beirlant \textit{et al.}(2018)]{BMV2018}Beirlant, J., Maribe, G.,
Verster, A., 2018. Penalized bias reduction in extreme value estimation for
censored Pareto-type data, and long-tailed insurance applications. Insurance
Math. Econom. \textbf{78}, 114-122.

\bibitem[Beirlant \textit{et al.}(2019)]{BWW2019}Beirlant, J., Worms, J.,
Worms, R., 2019. Estimation of the extreme value index in a censorship
framework: asymptotic and finite sample behavior. J. Statist. Plann. Inference
\textbf{202}, 31--56.

\bibitem[Bladt \textit{et al.}(2025)]{BGY2025}Bladt, M., Goegebeur, Y. and
Guillou, A. (2025). Asymptotically unbiased estimator of the extreme value
index under random censoring, https://hal.science/hal-04786783v2

\bibitem[Bladt and Rodionov(2025)]{BR2025}Bladt, M. and Rodionov, I. (2025).
Censored extreme value estimation, https://arxiv.org/abs/2312.10499.

\bibitem[Brahimi \textit{et} \textit{al.}(2015)]{BMN-2015}Brahimi, B.,
Meraghni, D., Necir, A., 2015. Approximations to the tail index estimator of a
heavy-tailed distribution under random censoring and application. Math.
Methods Statist. \textbf{24}, 266--279.

\bibitem[Caeiro \textit{et al.}(2005)]{Caeiro2005}Caeiro, F., Gomes, M. I.,
Pestana, D., 2005. Direct reduction of bias of the classical Hill estimator.
REVSTAT \textbf{3}, no. 2, 113--136

\bibitem[Colosimo \textit{et al.}(2002)]{Colosimo2002}Colosimo, E., Ferreira,
F. V., Oliveira, M., Sousa, C., 2002. Empirical comparisons between
Kaplan-Meier and Nelson-Aalen survival function estimators. J. Stat. Comput.
Simul. \textbf{72}, 299--308.

\bibitem[Cs\"{o}rg\H{o} \textit{et al.}(1986)]{CSCSHM86}Cs\"{o}rg\H{o}, M.,
Cs\"{o}rg\H{o}, S., Horv\'{a}th, L., Mason, D. M., 1986. Weighted empirical
and quantile processes. Ann. Probab. \textbf{14}, no. 1, 31--85.

\bibitem[Deheuvels and Einmahl(1996)]{Deheuvels1996}Deheuvels, P., Einmahl, J.
H. J, 1996. On the strong limiting behavior of local functionals of empirical
processes based upon censored data. Ann. Probab. \textbf{24}, no. 1, 504--525

\bibitem[Denuit \textit{et al.}(2006)]{DPV06}Denuit, M., Purcaru, O.,
Keilegom, I. V., 2006. Bivariate Archimedean copula models for censored data
in non-life insurance. J. Actuar. Pract. \textbf{13}, 5--32.

\bibitem[Einmahl and Koning(1992)]{Einmahl92}Einmahl, J. H. J., Koning, A. J.,
1992 Limit theorems for a general weighted process under random censoring.
Canad. J. Statist. \textbf{20}, no. 1, 77--89.

\bibitem[Einmahl \textit{et al.}(2006)]{EHL2006}Einmahl, J.H.J., de Haan, L.,
Li, D., 2006. Weighted approximations of tail copula processes with
application to testing the bivariate extreme value condition. Ann. Statist.
\textbf{34}, 1987--2014.

\bibitem[Einmahl \textit{et al.}(2008)]{EnFG08}Einmahl, J.H.J.,
Fils-Villetard, A., Guillou, A., 2008. Statistics of extremes under random
censoring. Bernoulli \textbf{14}, 207--227.

\bibitem[Frees and Valdez(1998)]{FV98}Frees, E. W., Valdez, E. A., 1998.
Understanding relationships using copulas. N. Am. Actuar. J. \textbf{2}, 1--25.

\bibitem[de Haan and Ferreira(2006)]{deHF06}de Haan, L., Ferreira, A., 2006.
Extreme Value Theory: An Introduction. Springer.

\bibitem[Hall(1982)]{Hall82}Hall, P., 1982. On some simple estimates of an
exponent of regular variation. J. Roy. Statist. Soc. Ser. B \textbf{44}, 37--42.

\bibitem[Hill(1975)]{Hill75}Hill, B.M., 1975. A simple general approach to
inference about the tail of a distribution. Ann. Statist. \textbf{3}, 1163-1174.

\bibitem[Hua and Joe(2011)]{HJ2011}Hua, L., and Joe, H., 2011. Second Order
Regular Variation and Conditional Tail Expectation of Multiple Risks.
Insurance Math. Econom. \textbf{49}, 537--546.

\bibitem[Kaplan and Meier(1958)]{KM58}Kaplan, E. L., Meier, P., 1958.
Nonparametric estimation from incomplete observations. J. Amer. Statist.
Assoc. \textbf{53}, 457--481.

\bibitem[Klugman and Parsa(1999)]{Klugman99}Klugman, S. A., Parsa, R., 1999.
Fitting bivariate loss distributions with copulas. Insurance Math. Econom.
\textbf{24}, 139--148.

\bibitem[Meraghni \textit{et al.}(2025)]{MNS2025}Meraghni, D., Necir, A.,
Soltane, L., 2025. Nelson-Aalen tail product-limit process and extreme value
index estimation under random censorship.\textit{ }Sankhya A\textit{. }https://doi.org/10.1007/s13171-025-00384-y.

\bibitem[Ndao et al.(2016)]{Ndao16}Ndao, P., Diop, A., Dupuy, J. F., 2016.
Nonparametric estimation of the conditional extreme-value index with random
covariates and censoring. J. Stat. Plan. Inference \textbf{168}, 20--37.

\bibitem[Nelson(1972)]{Nelson1972}Nelson, W., 1972. A short life test for
comparing a sample with previous accelerated test results. Technometrics
\textbf{14}, 175--185.

\bibitem[Neves and Fraga(2004)]{NF2004}Neves, C., Fraga, A. M. I., 2004. Reiss
and Thomas' automatic selection of the number of extremes. Comput. Statist.
Data Anal. \textbf{47}, no. 4, 689--704

\bibitem[Reiss and Thomas(1997)]{ReTo97}Reiss, R. D., Thomas, M., 2007.
Statistical analysis of extreme values: with applications to insurance,
finance, hydrology and other fields (3rd ed.). \textit{Birkh\"{a}user.}

\bibitem[Ripley and Solomon(1994)]{RS-94}Ripley, B. D., Solomon, P. J., 1994.
A note on Australian AIDS survival. \textit{University of Adelaide, Department
of Statistics, Research Report 94/3.}

\bibitem[Shorack and Wellner(1986)]{SW86}Shorack, G.R., Wellner, J.A., 1986.
Empirical Processes with Applications to Statistics. Wiley.

\bibitem[Venables and Ripley(2002)]{VR-02}Venables, W. N., Ripley, B. D.,
2002. Modern Applied Statistics with S (4th ed.). \textit{Springer.}

\bibitem[R. Worms and J. Worms(2014)]{WW2014}Worms, J., Worms, R., 2014. New
estimators of the extreme value index under random right censoring, for
heavy-tailed distributions. Extremes \textbf{17}, 337--358.
\end{thebibliography}

\section{\textbf{Appendix A: Auxiliary Results}}

\noindent In this section, we collect several auxiliary results that form the
technical foundation of the paper. For the sake of readability, only the
statements are presented here, while all detailed proofs are deferred to the
Supplementary Material.\smallskip

\noindent The results are organized according to their role in the analysis.
We first establish convergence properties and uniform bounds for the
Nelson--Aalen tail ratio and its nonlinear transformations. These preliminary
lemmas ensure the stability of the weighted structure and provide the
analytical tools required for handling powers of the empirical tail process.
We then derive structural representations of the estimator, which allow us to
express it as a functional of the underlying tail process. Next, we establish
analytical and probabilistic bounds for the stochastic components involved,
before controlling the boundary contribution induced by truncation. Finally,
we obtain a linearization of the estimator and identify the Gaussian limit
that governs its asymptotic behavior.\smallskip

\noindent Lemmas~\ref{lemma:weighted-power-limit} and
\ref{lemma:power-convergence-alpha} establish the convergence of powers of the
Nelson--Aalen tail ratio, both for fixed exponents and for random exponents
arising from the estimation of $p$. These results ensure that the nonlinear
transformations induced by the weighting scheme are asymptotically well
behaved.\smallskip

\noindent Lemma~\ref{lem:domination-Qn} provides uniform bounds for the tail
ratio and its powers over the truncated domain. This result is crucial for
establishing uniform integrability and for controlling the behavior of the
integrand in the functional representation of the estimator.\smallskip

\noindent Proposition~\ref{prop:app-tailratio} provides a uniform
approximation of the tail ratio under the regular variation assumption. This
result serves as a key ingredient in Proposition~\ref{prop:app-identity},
where it is used to derive an integral characterization of the extreme value
index adapted to the truncated setting.\smallskip

\noindent Proposition~\ref{prop:app-identity} establishes an integral
characterization of the extreme value index under truncation. This result
extends the classical identity to the weighted framework and plays a central
role in linking the estimator to its asymptotic target.\smallskip

\noindent Proposition~\ref{prop:app-functional} provides an exact functional
representation of the weighted and truncated Nelson--Aalen estimator. This
result is fundamental, as it expresses the estimator in terms of the empirical
tail process and serves as the starting point for all subsequent asymptotic
analyses.\smallskip

\noindent Proposition~\ref{prop:app-stochastic} establishes uniform stochastic
bounds for the processes $J_{n}(x)$ and $L_{n}(x)$ involved in the Gaussian
approximation of the weighted Nelson--Aalen tail process. These bounds ensure
uniform control of the stochastic fluctuations over the truncated
domain.\smallskip

\noindent Lemma~\ref{lem:boundary-truncation} controls the contribution of the
boundary term induced by truncation and shows that it is asymptotically
negligible under suitable conditions on the truncation sequence. This result
ensures that the truncation does not affect the asymptotic behavior of the
estimator.\smallskip

\noindent Proposition~\ref{prop:linearization} establishes a linearization of
the estimator at the $\sqrt{k}$ scale. It shows that the estimation error can
be decomposed into a linear functional of the weighted and truncated tail
process $D_{n,a}^{(\beta)}$, a deterministic bias term, and a negligible
remainder. This representation plays a central role in both the consistency
and asymptotic normality results.\smallskip

\noindent Proposition~\ref{prop:app-gaussian} identifies the Gaussian
component $\mathcal{N}_{n}(\beta)$ arising from the linearization and provides
its explicit variance. This result is crucial for deriving the asymptotic
normality of the estimator.\smallskip

\noindent Together, these auxiliary results provide a complete analytical and
probabilistic framework for establishing the consistency, stochastic
expansions, and asymptotic normality of the proposed estimators. In
particular, they highlight how the combination of truncation and weighting
restores a regular asymptotic structure over the full censoring range $0<p<1$.

\begin{proposition}
\label{prop:app-tailratio} Assume that the survival function $\overline{F}$
satisfies the first-order regular variation condition \eqref{RVF}. Then, for
every $\varepsilon>0$, there exists $t_{0}=t_{0}(\varepsilon)>0$ such that,
for all $t\geq t_{0}$ and all $y\geq1$ satisfying $ty\geq t_{0}$,
\[
\left|  \frac{\overline{F}(ty)}{\overline{F}(t)} - y^{-1/\gamma_{1}} \right|
\leq\varepsilon\max\!\left(  y^{-1/\gamma_{1}+\varepsilon}, y^{-1/\gamma
_{1}-\varepsilon} \right)  .
\]

\end{proposition}

\begin{proof}
This follows directly from the Potter-type bounds for regularly varying
functions; see Proposition~B.1.10 in \cite{deHF06}.
\end{proof}

\begin{proposition}
\label{prop:app-identity} Assume that the survival function $\overline{F}$
satisfies the first-order regular variation condition \eqref{RVF}. Let $u_{t}$
be a truncation level such that $u_{t}/t \to\infty$ as $t \to\infty$. Then,
for every $\beta>0$,
\begin{equation}
\frac{\beta}{p} \int_{1}^{u_{t}/t} y^{-1} \left(  \frac{\overline{F}%
(ty)}{\overline{F}(t)} \right)  ^{\beta/p} \,dy \longrightarrow\gamma_{1},
\qquad t \to\infty. \label{eq:app-identity-y}%
\end{equation}
Equivalently,
\begin{equation}
\left(  \frac{\beta}{p}\right)  ^{2} \int_{t}^{u_{t}} \left(  \frac
{\overline{F}(y)}{\overline{F}(t)} \right)  ^{\beta/p-1} \log\!\left(
\frac{y}{t}\right)  \,d\!\left(  \frac{F(y)}{\overline{F}(t)} \right)
\longrightarrow\gamma_{1}, \qquad t \to\infty. \label{eq:app-identity-t}%
\end{equation}

\end{proposition}

\begin{proof}
Set $a:=\beta/p.$ By Proposition~\ref{prop:app-tailratio}, for every
$\varepsilon>0$, there exists $t_{0}>0$ such that, for all $t\geq t_{0}$ and
$ty\geq t_{0}$,
\[
\left\vert \frac{\overline{F}(ty)}{\overline{F}(t)}-y^{-1/\gamma_{1}%
}\right\vert \leq\varepsilon\max\left(  y^{-1/\gamma_{1}+\varepsilon
},\,y^{-1/\gamma_{1}-\varepsilon}\right)  .
\]
In particular, for $y\geq1$ and $t\geq t_{0}$, this implies
\[
\frac{\overline{F}(ty)}{\overline{F}(t)}\leq y^{-1/\gamma_{1}}+\varepsilon
y^{-1/\gamma_{1}+\varepsilon}\leq C\,y^{-1/\gamma_{1}+\varepsilon},
\]
for some constant $C>0$. Hence,
\[
y^{-1}\left(  \frac{\overline{F}(ty)}{\overline{F}(t)}\right)  ^{a}\leq
C\,y^{-1-a/\gamma_{1}+a\varepsilon}.
\]
Choosing $\varepsilon>0$ sufficiently small so that
\[
\frac{a}{\gamma_{1}}-a\varepsilon>0,
\]
the exponent $-1-a/\gamma_{1}+a\varepsilon$ is strictly less than $-1$, and
thus the function $y\mapsto y^{-1-a/\gamma_{1}+a\varepsilon}$ is integrable on
$[1,\infty)$. On the other hand, by regular variation,
\[
\left(  \frac{\overline{F}(ty)}{\overline{F}(t)}\right)  ^{a}\longrightarrow
y^{-a/\gamma_{1}},\qquad t\rightarrow\infty,
\]
for each fixed $y\geq1$. Using a standard truncation argument together with
dominated convergence on compact intervals, and letting the truncation level
tend to infinity, we obtain
\[
a\int_{1}^{u_{t}/t}y^{-1}\left(  \frac{\overline{F}(ty)}{\overline{F}%
(t)}\right)  ^{a}dy\longrightarrow a\int_{1}^{\infty}y^{-1-a/\gamma_{1}}\,dy.
\]
Since
\[
a\int_{1}^{\infty}y^{-1-a/\gamma_{1}}\,dy=a\cdot\frac{\gamma_{1}}{a}%
=\gamma_{1},
\]
this proves \eqref{eq:app-identity-y}.\smallskip
\end{proof}

\begin{proposition}
\label{prop:app-stochastic} Let $1/4<\eta<1/2$ and let $\varepsilon>0$ be
sufficiently small. Then, uniformly for $1\leq x\leq T_{n,a}$, the following
assertions hold:
\[
(i)\qquad J_{n}(x)=O_{\mathbf{P}}\!\left(  x^{(2\eta-p)/\gamma+\varepsilon
}\right)  ,
\]%
\[
(ii)\qquad\sqrt{k}\,\left(  \frac{\overline{F}_{n}^{(\mathrm{NA})}%
(Z_{n-k:n}x)}{\overline{F}_{n}^{(\mathrm{NA})}(Z_{n-k:n})}-\frac{\overline
{F}(Z_{n-k:n}x)}{\overline{F}(Z_{n-k:n})}\right)  =O_{\mathbf{P}}\!\left(
x^{(2\eta-p)/\gamma+\varepsilon}\right)  ,
\]
and
\[
(iii)\qquad\frac{\overline{F}_{n}^{(\mathrm{NA})}(Z_{n-k:n}x)}{\overline
{F}_{n}^{(\mathrm{NA})}(Z_{n-k:n})}=x^{-p/\gamma}\left(  1+O_{\mathbf{P}%
}\!\left(  k^{-1/2}x^{2\eta/\gamma+\varepsilon}\right)  +o_{\mathbf{P}%
}(1)\right)  ,
\]

\end{proposition}

\begin{proof}
We first recall the standard Brownian bound: for every $0<\nu<1/2$,
\[
\sup_{0<s\leq1}\frac{|W_{n}(s)|}{s^{\nu}}=O_{\mathbf{P}}(1).
\]
Choosing $\nu=1-2\eta$, we obtain
\[
|W_{n}(s)|=O_{\mathbf{P}}(s^{1-2\eta}),\qquad\text{uniformly for }0<s\leq1.
\]
The same bound holds for the processes $\mathbf{W}_{n,1}$ and $\mathbf{W}%
_{n,2}$.\smallskip

\noindent We first consider
\[
J_{n1}(x)=x^{1/\gamma_{2}}\mathbf{W}_{n,1}(x^{-1/\gamma})-x^{-1/\gamma_{1}%
}\mathbf{W}_{n,1}(1).
\]
Since
\[
\mathbf{W}_{n,1}(x^{-1/\gamma})=O_{\mathbf{P}}\!\left(  x^{-(1-2\eta)/\gamma
}\right)  ,
\]
uniformly for $1\leq x\leq T_{n,a}$, and using
\[
\frac{1}{\gamma_{2}}=\frac{1-p}{\gamma},\qquad\frac{1}{\gamma_{1}}=\frac
{p}{\gamma},
\]
we obtain
\[
x^{1/\gamma_{2}}\mathbf{W}_{n,1}(x^{-1/\gamma})=O_{\mathbf{P}}\!\left(
x^{(2\eta-p)/\gamma}\right)  .
\]
Moreover,
\[
x^{-1/\gamma_{1}}\mathbf{W}_{n,1}(1)=O_{\mathbf{P}}\!\left(  x^{-p/\gamma
}\right)  =O_{\mathbf{P}}\!\left(  x^{(2\eta-p)/\gamma}\right)  ,
\]
since $x\geq1$. Hence
\[
J_{n1}(x)=O_{\mathbf{P}}\!\left(  x^{(2\eta-p)/\gamma}\right)  .
\]
Next, consider
\[
J_{n2}(x)=\gamma^{-1}x^{-1/\gamma_{1}}\int_{1}^{x}u^{1/\gamma-1}\left(
p\mathbf{W}_{n,2}(u^{-1/\gamma})-q\mathbf{W}_{n,1}(u^{-1/\gamma})\right)
\,du.
\]
Using the same Brownian bound,
\[
\mathbf{W}_{n,i}(u^{-1/\gamma})=O_{\mathbf{P}}\!\left(  u^{-(1-2\eta)/\gamma
}\right)  ,\qquad i=1,2,
\]
uniformly for $1\leq u\leq x$, so the integrand is of order
\[
u^{1/\gamma-1}u^{-(1-2\eta)/\gamma}=u^{2\eta/\gamma-1}.
\]
It follows that
\[
\int_{1}^{x}u^{2\eta/\gamma-1}\,du=O\!\left(  x^{2\eta/\gamma}\right)  .
\]
Multiplying by $x^{-1/\gamma_{1}}=x^{-p/\gamma}$, we obtain
\[
J_{n2}(x)=O_{\mathbf{P}}\!\left(  x^{(2\eta-p)/\gamma}\right)  .
\]
Thus,
\[
J_{n}(x)=J_{n1}(x)+J_{n2}(x)=O_{\mathbf{P}}\!\left(  x^{(2\eta-p)/\gamma
}\right)  ,
\]
which yields assertion (i).\smallskip

\noindent Assertion (ii) follows from the Gaussian approximation established
in Theorem~\ref{theorem1}, which gives
\[
\sqrt{k}\,L_{n}(x)=J_{n}(x)+o_{\mathbf{P}}\!\left(  x^{(2\eta-p)/\gamma
+\varepsilon}\right)  ,
\]
uniformly for $1\leq x\leq T_{n,a}$. Combining this with (i) yields
\[
\sqrt{k}\left(  \frac{\overline{F}_{n}^{(\mathrm{NA})}(Z_{n-k:n}x)}%
{\overline{F}_{n}^{(\mathrm{NA})}(Z_{n-k:n})}-\frac{\overline{F}(Z_{n-k:n}%
x)}{\overline{F}(Z_{n-k:n})}\right)  =O_{\mathbf{P}}\!\left(  x^{(2\eta
-p)/\gamma+\varepsilon}\right)  .
\]
We write
\[
\frac{\overline{F}_{n}^{(\mathrm{NA})}(Z_{n-k:n}x)}{\overline{F}%
_{n}^{(\mathrm{NA})}(Z_{n-k:n})}=\frac{\overline{F}(Z_{n-k:n}x)}{\overline
{F}(Z_{n-k:n})}+L_{n}(x),
\]
where
\[
L_{n}(x):=\frac{\overline{F}_{n}^{(\mathrm{NA})}(Z_{n-k:n}x)}{\overline{F}%
_{n}^{(\mathrm{NA})}(Z_{n-k:n})}-\frac{\overline{F}(Z_{n-k:n}x)}{\overline
{F}(Z_{n-k:n})}.
\]
Since $Z_{n-k:n}\rightarrow\infty$ in probability,
Proposition~\ref{prop:app-tailratio} yields
\[
\frac{\overline{F}(Z_{n-k:n}x)}{\overline{F}(Z_{n-k:n})}=x^{-1/\gamma_{1}%
}\left(  1+o_{\mathbf{P}}(1)\right)  .
\]
Moreover, by assertion (ii),
\[
L_{n}(x)=O_{\mathbf{P}}\!\left(  k^{-1/2}x^{(2\eta-p)/\gamma+\varepsilon
}\right)  .
\]
Dividing by $x^{-p/\gamma}$ gives
\[
\frac{L_{n}(x)}{x^{-p/\gamma}}=O_{\mathbf{P}}\!\left(  k^{-1/2}x^{2\eta
/\gamma+\varepsilon}\right)  .
\]
Hence,
\[
\frac{\overline{F}_{n}^{(\mathrm{NA})}(Z_{n-k:n}x)}{\overline{F}%
_{n}^{(\mathrm{NA})}(Z_{n-k:n})}=x^{-p/\gamma}\left[  1+O_{\mathbf{P}%
}\!\left(  k^{-1/2}x^{2\eta/\gamma+\varepsilon}\right)  +o_{\mathbf{P}%
}(1)\right]  ,
\]
uniformly for $1\leq x\leq T_{n,a}$. This proves (iii).
\end{proof}

\begin{proposition}
\label{prop:app-gaussian} Let $\beta>1$ and $0<p<1$. Define
\[
\mathcal{N}_{n}(\beta) = \frac{\beta(p+\beta-1)}{p} \int_{0}^{1}s^{\beta
-2}\mathbf{W}_{n,1}(s)\,ds + \beta\int_{0}^{1}s^{\beta-2}\mathbf{W}%
_{n,2}(s)\,ds - \frac{\beta}{p}\mathbf{W}_{n,1}(1).
\]
Then $\mathcal{N}_{n}(\beta)$ is a centered Gaussian random variable and
\[
\operatorname{Var}\!\left(  \mathcal{N}_{n}(\beta)\right)  = \frac{1}{p}%
\frac{\beta^{2}}{2\beta-1}.
\]

\end{proposition}

\begin{proof}
The random variable $\mathcal{N}_{n}(\beta)$ is a linear functional of the
centered Gaussian processes $\mathbf{W}_{n,1}$ and $\mathbf{W}_{n,2}$. Hence
it is centered Gaussian. It remains to compute its variance.\smallskip

\noindent Since $\mathbf{W}_{n,1}$ and $\mathbf{W}_{n,2}$ are independent and
satisfy
\[
\mathbb{E}[\mathbf{W}_{n,1}(s)\mathbf{W}_{n,1}(t)]=p\min(s,t),\qquad
\mathbb{E}[\mathbf{W}_{n,2}(s)\mathbf{W}_{n,2}(t)]=q\min(s,t),
\]
we have, for $\beta>1$,
\[
\operatorname{Var}\left(  \int_{0}^{1}s^{\beta-2}\mathbf{W}_{n,i}%
(s)\,ds\right)  =c_{i}\frac{2}{\beta(2\beta-1)},
\]
where $c_{1}=p$ and $c_{2}=q$.\smallskip

\noindent Moreover,
\[
\operatorname{Cov}\left(  \int_{0}^{1}s^{\beta-2}\mathbf{W}_{n,1}%
(s)\,ds,\mathbf{W}_{n,1}(1)\right)  =p\int_{0}^{1}s^{\beta-1}\,ds=\frac
{p}{\beta}.
\]
Therefore,
\[
\begin{aligned}
		\operatorname{Var}\!\left(\mathcal{N}_{n}(\beta)\right)
		&=
		\left(\frac{\beta(p+\beta-1)}{p}\right)^{2}
		p\frac{2}{\beta(2\beta-1)}
		+
		\beta^{2}q\frac{2}{\beta(2\beta-1)}
		\\
		&\quad
		+
		\left(\frac{\beta}{p}\right)^{2}p
		-
		2\frac{\beta(p+\beta-1)}{p}\frac{\beta}{p}\frac{p}{\beta}.
	\end{aligned}
\]
A direct simplification gives
\[
\operatorname{Var}\!\left(  \mathcal{N}_{n}(\beta)\right)  =\frac{1}{p}%
\frac{\beta^{2}}{2\beta-1}.
\]
This completes the proof.
\end{proof}

\begin{proposition}
\label{prop:app-functional} Set
\[
Z_{k}:=Z_{n-k:n}, \qquad T_{n,a}:=\frac{u_{n}}{Z_{k}}, \qquad\alpha_{n}%
:=\frac{\beta}{\widehat{p}_{k}},
\]
and, for $1\leq x\leq T_{n,a}$, define
\[
Q_{n}(x) := \frac{\overline{F}_{n}^{(\mathrm{NA})}(xZ_{k})} {\overline{F}%
_{n}^{(\mathrm{NA})}(Z_{k})}, \qquad R_{n,a}(x) := \frac{\overline{F}%
_{n}^{(\mathrm{NA},a)}(xZ_{k})} {\overline{F}_{n}^{(\mathrm{NA})}(Z_{k})},
\]
and
\[
c_{n} := \frac{\overline{F}_{n}^{(\mathrm{NA})}(u_{n})} {\overline{F}%
_{n}^{(\mathrm{NA})}(Z_{k})}.
\]
Then, for every $\beta>1$,
\begin{equation}
\widehat{\gamma}_{1,k}^{(\mathrm{NA,tr})}(\beta) = \alpha_{n} \int%
_{1}^{T_{n,a}} x^{-1} \left(  Q_{n}(x)\right)  ^{\alpha_{n}} \,dx - \alpha
_{n}(\log T_{n,a})\,c_{n}^{\alpha_{n}}. \label{eq:app-functional-Q}%
\end{equation}
Moreover, for $1\leq x\leq T_{n,a}$,
\[
Q_{n}(x)=R_{n,a}(x)+c_{n},
\]
so that
\begin{equation}
\widehat{\gamma}_{1,k}^{(\mathrm{NA,tr})}(\beta) = \alpha_{n} \int%
_{1}^{T_{n,a}} x^{-1} \left(  R_{n,a}(x)+c_{n}\right)  ^{\alpha_{n}} \,dx -
\alpha_{n}(\log T_{n,a})\,c_{n}^{\alpha_{n}}. \label{eq:app-functional-R}%
\end{equation}
Finally, using
\[
D_{n,a}^{(\beta)}(x) = \sqrt{k}\,x^{-\beta/\gamma} \left(  R_{n,a}%
(x)-x^{-1/\gamma_{1}} \right)  , \qquad1\leq x\leq T_{n,a},
\]
that is,
\[
R_{n,a}(x) = x^{-1/\gamma_{1}} + \frac{1}{\sqrt{k}}x^{\beta/\gamma}%
D_{n,a}^{(\beta)}(x),
\]
we obtain the exact functional representation
\begin{equation}
\widehat{\gamma}_{1,k}^{(\mathrm{NA,tr})}(\beta) = \alpha_{n} \int%
_{1}^{T_{n,a}} x^{-1} \left(  x^{-1/\gamma_{1}} + c_{n} + \frac{1}{\sqrt{k}%
}x^{\beta/\gamma}D_{n,a}^{(\beta)}(x) \right)  ^{\alpha_{n}} dx - \alpha
_{n}(\log T_{n,a})\,c_{n}^{\alpha_{n}}. \label{eq:app-functional-D}%
\end{equation}

\end{proposition}

\begin{proof}
Recall that the weighted and truncated Nelson--Aalen estimator is defined by
\[
\widehat{\gamma}_{1,k}^{(\mathrm{NA,tr})}(\beta) = \alpha_{n}^{2} \int_{Z_{k}%
}^{u_{n}} \left(  \frac{\overline{F}_{n}^{(\mathrm{NA})}(z)} {\overline{F}%
_{n}^{(\mathrm{NA})}(Z_{k})} \right)  ^{\alpha_{n}-1} \log\!\left(  \frac
{z}{Z_{k}}\right)  \,d\!\left(  \frac{F_{n}^{(\mathrm{NA})}(z)} {\overline
{F}_{n}^{(\mathrm{NA})}(Z_{k})} \right)  ,
\]
where
\[
\alpha_{n}=\frac{\beta}{\widehat{p}_{k}}.
\]
We make the change of variables
\[
z=xZ_{k},\qquad1\leq x\leq T_{n,a}:=\frac{u_{n}}{Z_{k}}.
\]
Then
\[
\log\!\left(  \frac{z}{Z_{k}}\right)  =\log x, \qquad Q_{n}(x)= \frac
{\overline{F}_{n}^{(\mathrm{NA})}(xZ_{k})} {\overline{F}_{n}^{(\mathrm{NA}%
)}(Z_{k})}.
\]
Since
\[
Q_{n}(x) = 1- \frac{F_{n}^{(\mathrm{NA})}(xZ_{k})} {\overline{F}%
_{n}^{(\mathrm{NA})}(Z_{k})},
\]
we have
\[
d\!\left(  \frac{F_{n}^{(\mathrm{NA})}(xZ_{k})} {\overline{F}_{n}%
^{(\mathrm{NA})}(Z_{k})} \right)  = -\,dQ_{n}(x).
\]
Therefore,
\[
\widehat{\gamma}_{1,k}^{(\mathrm{NA,tr})}(\beta) = -\alpha_{n}^{2} \int%
_{1}^{T_{n,a}} Q_{n}(x)^{\alpha_{n}-1}\log x\,dQ_{n}(x).
\]
Since
\[
d\!\left(  Q_{n}(x)^{\alpha_{n}}\right)  = \alpha_{n}Q_{n}(x)^{\alpha_{n}%
-1}\,dQ_{n}(x),
\]
we obtain
\[
\widehat{\gamma}_{1,k}^{(\mathrm{NA,tr})}(\beta) = -\alpha_{n} \int%
_{1}^{T_{n,a}} \log x\,d\!\left(  Q_{n}(x)^{\alpha_{n}}\right)  .
\]
Since $Q_{n}$ is monotone, the integration by parts formula applies. Hence,
\[
-\alpha_{n} \int_{1}^{T_{n,a}} \log x\,d\!\left(  Q_{n}(x)^{\alpha_{n}%
}\right)  = -\alpha_{n} \left[  \log x\,Q_{n}(x)^{\alpha_{n}} \right]
_{1}^{T_{n,a}} + \alpha_{n} \int_{1}^{T_{n,a}} x^{-1}Q_{n}(x)^{\alpha_{n}%
}\,dx.
\]
Since $\log1=0$ and
\[
Q_{n}(T_{n,a}) = \frac{\overline{F}_{n}^{(\mathrm{NA})}(u_{n})} {\overline
{F}_{n}^{(\mathrm{NA})}(Z_{k})} = c_{n},
\]
we get
\[
\widehat{\gamma}_{1,k}^{(\mathrm{NA,tr})}(\beta) = \alpha_{n} \int%
_{1}^{T_{n,a}} x^{-1}\left(  Q_{n}(x)\right)  ^{\alpha_{n}}\,dx - \alpha
_{n}(\log T_{n,a})c_{n}^{\alpha_{n}}.
\]
This proves \eqref{eq:app-functional-Q}.\smallskip

\noindent Next, by definition of the truncated tail,
\[
\overline{F}_{n}^{(\mathrm{NA},a)}(xZ_{k}) = \overline{F}_{n}^{(\mathrm{NA}%
)}(xZ_{k}) - \overline{F}_{n}^{(\mathrm{NA})}(u_{n}), \qquad1\leq x\leq
T_{n,a}.
\]
Dividing both sides by $\overline{F}_{n}^{(\mathrm{NA})}(Z_{k})$ gives
\[
R_{n,a}(x)=Q_{n}(x)-c_{n},
\]
or equivalently,
\[
Q_{n}(x)=R_{n,a}(x)+c_{n}.
\]
Substituting this identity into \eqref{eq:app-functional-Q} yields
\[
\widehat{\gamma}_{1,k}^{(\mathrm{NA,tr})}(\beta) = \alpha_{n} \int%
_{1}^{T_{n,a}} x^{-1}\left(  R_{n,a}(x)+c_{n}\right)  ^{\alpha_{n}}\,dx -
\alpha_{n}(\log T_{n,a})c_{n}^{\alpha_{n}},
\]
which proves \eqref{eq:app-functional-R}.\smallskip

\noindent Finally, from the definition of the weighted truncated process,
\[
D_{n,a}^{(\beta)}(x) = \sqrt{k}\,x^{-\beta/\gamma} \left(  R_{n,a}%
(x)-x^{-1/\gamma_{1}}\right)  ,
\]
we obtain
\[
R_{n,a}(x) = x^{-1/\gamma_{1}} + \frac{1}{\sqrt{k}}x^{\beta/\gamma}%
D_{n,a}^{(\beta)}(x).
\]
Substituting this identity into \eqref{eq:app-functional-R} gives
\[
\widehat{\gamma}_{1,k}^{(\mathrm{NA,tr})}(\beta) = \alpha_{n} \int%
_{1}^{T_{n,a}} x^{-1} \left(  x^{-1/\gamma_{1}} + c_{n} + \frac{1}{\sqrt{k}%
}x^{\beta/\gamma}D_{n,a}^{(\beta)}(x) \right)  ^{\alpha_{n}}dx - \alpha
_{n}(\log T_{n,a})c_{n}^{\alpha_{n}}.
\]
This proves \eqref{eq:app-functional-D} and completes the proof.
\end{proof}

\begin{lemma}
\label{lem:boundary-truncation} Assume that $(a_{n})$ satisfies
\[
a_{n}\to0,\qquad na_{n}\to\infty,\qquad a_{n}=o(k/n),
\]
and, in addition,
\[
\sqrt{k} \left(  \frac{na_{n}}{k}\right)  ^{\beta} \log\!\left(  \frac
{k}{na_{n}}\right)  \to0.
\]
Then
\[
\sqrt{k}\,\alpha_{n}(\log T_{n,a})c_{n}^{\alpha_{n}} = o_{\mathbf{P}}(1).
\]

\end{lemma}

\begin{proof}
Recall that
\[
c_{n}=\frac{\overline{F}_{n}^{(\mathrm{NA})}(u_{n})}{\overline{F}%
_{n}^{(\mathrm{NA})}(Z_{k})},\qquad T_{n,a}=\frac{u_{n}}{Z_{k}}.
\]
By the uniform consistency of the Nelson--Aalen tail ratio on the truncated
domain, together with Potter-type bounds, we have
\[
c_{n}=O_{\mathbf{P}}\!\left(  \frac{\overline{F}(u_{n})}{\overline{F}(Z_{k}%
)}\right)  .
\]
Since $u_{n}$ corresponds to a tail level of order $a_{n}$, while $Z_{k}$
corresponds to a tail level of order $k/n$, regular variation yields
\[
\frac{\overline{F}(u_{n})}{\overline{F}(Z_{k})}=O\!\left[  \left(
\frac{na_{n}}{k}\right)  ^{p}\right]  .
\]
Consequently,
\[
c_{n}=O_{\mathbf{P}}\!\left[  \left(  \frac{na_{n}}{k}\right)  ^{p}\right]  .
\]
Moreover, since
\[
\alpha_{n}=\frac{\beta}{\widehat{p}_{k}}\qquad\text{and}\qquad\widehat{p}%
_{k}\overset{\mathbf{P}}{\longrightarrow}p,
\]
we have
\[
\alpha_{n}\overset{\mathbf{P}}{\longrightarrow}\frac{\beta}{p}.
\]
Therefore,
\[
c_{n}^{\alpha_{n}}=O_{\mathbf{P}}\!\left[  \left(  \frac{na_{n}}{k}\right)
^{\beta}\right]  .
\]
It remains to control the logarithmic factor. Since $u_{n}$ and $Z_{k}$
correspond respectively to the tail levels $a_{n}$ and $k/n$, regular
variation of the tail quantile implies
\[
T_{n,a}=\frac{u_{n}}{Z_{k}}=O_{\mathbf{P}}\!\left[  \left(  \frac{k}{na_{n}%
}\right)  ^{\gamma}\right]  .
\]
Hence,
\[
\log T_{n,a}=O_{\mathbf{P}}\!\left[  \log\!\left(  \frac{k}{na_{n}}\right)
\right]  .
\]
Combining the previous estimates gives
\[
\sqrt{k}\,\alpha_{n}(\log T_{n,a})c_{n}^{\alpha_{n}}=O_{\mathbf{P}}\!\left[
\sqrt{k}\left(  \frac{na_{n}}{k}\right)  ^{\beta}\log\!\left(  \frac{k}%
{na_{n}}\right)  \right]  .
\]
By the imposed rate condition, the right-hand side is $o_{\mathbf{P}}(1)$.
This proves the lemma.
\end{proof}

\begin{proposition}
[Linearization of the estimator]\label{prop:linearization} Assume that the
conditions of Theorem~\ref{theorem2} hold. Assume, in addition, that
\[
\sqrt{k} \left(  \frac{na_{n}}{k}\right)  ^{\beta} \log\!\left(  \frac
{k}{na_{n}}\right)  \to0.
\]
Then, for every $\beta>1$, we have
\[
\sqrt{k}\left(  \widehat{\gamma}_{1,k}^{(\mathrm{NA,tr})}(\beta) - \gamma_{1}
\right)  = \left(  \frac{\beta}{p}\right)  ^{2} \int_{1}^{T_{n,a}}
x^{-1+1/\gamma_{1}} D_{n,a}^{(\beta)}(x)\,dx + B_{n} + o_{\mathbf{P}}(1),
\]
where
\[
B_{n} := \sqrt{k} \left[  \frac{\beta}{p} \int_{1}^{T_{n,a}} x^{-1} \left(
\frac{\overline{F}(Z_{k}x)}{\overline{F}(Z_{k})} \right)  ^{\beta/p} \,dx -
\gamma_{1} \right]  .
\]

\end{proposition}

\begin{proof}
We start from the exact functional representation established in
Proposition~\ref{prop:app-functional}:
\[
\widehat{\gamma}_{1,k}^{(\mathrm{NA,tr})}(\beta)=\alpha_{n}\int_{1}^{T_{n,a}%
}x^{-1}Q_{n}(x)^{\alpha_{n}}\,dx-\alpha_{n}(\log T_{n,a})c_{n}^{\alpha_{n}},
\]
where
\[
\alpha_{n}:=\frac{\beta}{\widehat{p}_{k}}.
\]
Since $\widehat{p}_{k}\overset{\mathbf{P}}{\longrightarrow}p$, we have
\[
\alpha_{n}\overset{\mathbf{P}}{\longrightarrow}\frac{\beta}{p}.
\]
Set
\[
A_{n}(x):=\frac{\overline{F}(Z_{k}x)}{\overline{F}(Z_{k})}.
\]
We decompose
\[
\begin{aligned}
		\sqrt{k}\left(
		\widehat{\gamma}_{1,k}^{(\mathrm{NA,tr})}(\beta)-\gamma_{1}
		\right)
		&=
		\sqrt{k}
		\left[
		\alpha_n\int_{1}^{T_{n,a}}x^{-1}Q_n(x)^{\alpha_n}\,dx
		-
		\frac{\beta}{p}\int_{1}^{T_{n,a}}x^{-1}A_n(x)^{\beta/p}\,dx
		\right]  \\
		&\quad+
		\sqrt{k}
		\left[
		\frac{\beta}{p}\int_{1}^{T_{n,a}}x^{-1}A_n(x)^{\beta/p}\,dx
		-\gamma_1
		\right]  \\
		&\quad-
		\sqrt{k}\,\alpha_n(\log T_{n,a})c_n^{\alpha_n}.
	\end{aligned}
\]
The second term is exactly $B_{n}$. By Lemma~\ref{lem:boundary-truncation},
the last term is $o_{\mathbf{P}}(1)$.

It remains to linearize the first term. By the uniform stochastic expansion of
the Nelson--Aalen tail ratio,
\[
Q_{n}(x)=A_{n}(x)+\{Q_{n}(x)-A_{n}(x)\},
\]
and Taylor's formula applied to the map $u\mapsto u^{\beta/p}$ gives
\[
Q_{n}(x)^{\beta/p}=A_{n}(x)^{\beta/p}+\frac{\beta}{p}A_{n}(x)^{\beta
/p-1}\{Q_{n}(x)-A_{n}(x)\}+r_{n}(x),
\]
where
\[
\sqrt{k}\int_{1}^{T_{n,a}}x^{-1}|r_{n}(x)|\,dx=o_{\mathbf{P}}(1).
\]
Moreover, the replacement of $\alpha_{n}$ by $\beta/p$ in the preceding
integral contributes only $o_{\mathbf{P}}(1)$ after multiplication by
$\sqrt{k}$. Hence
\[
\begin{aligned}
		&\sqrt{k}
		\left[
		\alpha_n\int_{1}^{T_{n,a}}x^{-1}Q_n(x)^{\alpha_n}\,dx
		-
		\frac{\beta}{p}\int_{1}^{T_{n,a}}x^{-1}A_n(x)^{\beta/p}\,dx
		\right] \\
		&\qquad =
		\left(\frac{\beta}{p}\right)^{2}
		\int_{1}^{T_{n,a}}x^{-1}A_n(x)^{\beta/p-1}
		\sqrt{k}\{Q_n(x)-A_n(x)\}\,dx
		+
		o_{\mathbf P}(1).
	\end{aligned}
\]
From the definition of the weighted truncated process,
\[
D_{n,a}^{(\beta)}(x)=\sqrt{k}\,x^{-\beta/\gamma}\left(  R_{n,a}%
(x)-x^{-1/\gamma_{1}}\right)  ,
\]
and from the fact that the truncation correction is negligible on the scale
considered, we have
\[
\sqrt{k}\{Q_{n}(x)-A_{n}(x)\}=x^{\beta/\gamma}D_{n,a}^{(\beta)}%
(x)+o_{\mathbf{P}}(1)
\]
in the corresponding weighted integral sense. Furthermore, by regular
variation,
\[
A_{n}(x)^{\beta/p-1}=x^{-(\beta/p-1)/\gamma_{1}}\{1+o_{\mathbf{P}}(1)\},
\]
uniformly in the same weighted integral sense. Since $p=\gamma/\gamma_{1}$,
\[
-\frac{\beta/p-1}{\gamma_{1}}+\frac{\beta}{\gamma}=\frac{1}{\gamma_{1}}.
\]
Therefore,
\[
x^{-1}A_{n}(x)^{\beta/p-1}\sqrt{k}\{Q_{n}(x)-A_{n}(x)\}=x^{-1+1/\gamma_{1}%
}D_{n,a}^{(\beta)}(x)+o_{\mathbf{P}}(1)
\]
in the weighted integral sense. Consequently,
\[
\begin{aligned}
		&\sqrt{k}
		\left[
		\alpha_n\int_{1}^{T_{n,a}}x^{-1}Q_n(x)^{\alpha_n}\,dx
		-
		\frac{\beta}{p}\int_{1}^{T_{n,a}}x^{-1}A_n(x)^{\beta/p}\,dx
		\right] \\
		&\qquad =
		\left(\frac{\beta}{p}\right)^{2}
		\int_{1}^{T_{n,a}}x^{-1+1/\gamma_1}D_{n,a}^{(\beta)}(x)\,dx
		+
		o_{\mathbf P}(1).
	\end{aligned}
\]
Combining this with the definition of $B_{n}$ and
Lemma~\ref{lem:boundary-truncation} gives the desired result.
\end{proof}

\begin{lemma}
\label{lemma:weighted-power-limit} Let $\beta>1$ and $0<p<1$. Assume that the
survival function $\overline{F}$ satisfies the first-order regular variation
condition \eqref{RVF}. Then, for each fixed $x\geq1$,
\[
\left(  Q_{n}(x)\right)  ^{\beta/p} \overset{\mathbf{P}}{\longrightarrow}
x^{-\beta/\gamma}.
\]

\end{lemma}

\begin{proof}
Recall that
\[
Q_{n}(x) = \frac{\overline{F}_{n}^{(\mathrm{NA})}(Z_{k}x)} {\overline{F}%
_{n}^{(\mathrm{NA})}(Z_{k})}.
\]
By the consistency of the Nelson--Aalen tail ratio (see
Proposition~\ref{prop:app-stochastic}), we have, for each fixed $x\geq1$,
\[
Q_{n}(x) = \frac{\overline{F}(Z_{k}x)}{\overline{F}(Z_{k})} + o_{\mathbf{P}%
}(1).
\]
Since $\overline{F}$ is regularly varying with index $-1/\gamma_{1}$, we have
\[
\frac{\overline{F}(Z_{k}x)}{\overline{F}(Z_{k})} \overset{\mathbf{P}%
}{\longrightarrow} x^{-1/\gamma_{1}}.
\]
Therefore,
\[
Q_{n}(x) \overset{\mathbf{P}}{\longrightarrow} x^{-1/\gamma_{1}}.
\]
By continuity of the map $u\mapsto u^{\beta/p}$ on $(0,\infty)$, it follows
that
\[
\left(  Q_{n}(x)\right)  ^{\beta/p} \overset{\mathbf{P}}{\longrightarrow}
\left(  x^{-1/\gamma_{1}}\right)  ^{\beta/p}.
\]
Using the identity
\[
\frac{1}{\gamma_{1}}=\frac{p}{\gamma},
\]
we obtain
\[
\left(  x^{-1/\gamma_{1}}\right)  ^{\beta/p} = x^{-\beta/\gamma},
\]
which completes the proof.
\end{proof}

\begin{lemma}
\label{lemma:power-convergence-alpha} Let $\beta>1$ and $0<p<1$, and define
\[
\alpha_{n}:=\frac{\beta}{\widehat{p}_{k}}, \qquad\alpha_{n}\overset{\mathbf{P}%
}{\longrightarrow}\frac{\beta}{p}.
\]
Let $Q_{n}(x)$ denote the Nelson--Aalen tail ratio. Then, for each fixed
$x\geq1$,
\[
\left(  Q_{n}(x)\right)  ^{\alpha_{n}} \overset{\mathbf{P}}{\longrightarrow}
x^{-\beta/\gamma}.
\]

\end{lemma}

\begin{proof}
From Lemma~\ref{lemma:weighted-power-limit}, we have, for each fixed $x\geq
1$,
\[
Q_{n}(x)\overset{\mathbf{P}}{\longrightarrow} x^{-1/\gamma_{1}}.
\]
Moreover,
\[
\alpha_{n} \overset{\mathbf{P}}{\longrightarrow} \frac{\beta}{p}.
\]
Since $x^{-1/\gamma_{1}}>0$, we can apply the continuous mapping theorem to
the function
\[
(u,a)\mapsto u^{a},
\]
which is continuous on $(0,\infty)\times\mathbb{R}$. Hence,
\[
\left(  Q_{n}(x)\right)  ^{\alpha_{n}} \overset{\mathbf{P}}{\longrightarrow}
\left(  x^{-1/\gamma_{1}}\right)  ^{\beta/p}.
\]
Using the identity
\[
\frac{1}{\gamma_{1}}=\frac{p}{\gamma},
\]
we obtain
\[
\left(  x^{-1/\gamma_{1}}\right)  ^{\beta/p} = x^{-\beta/\gamma},
\]
which completes the proof.
\end{proof}

\begin{lemma}
\label{lem:domination-Qn} Let $\beta>1$ and $0<p<1$. Assume that the
conditions of Theorem~\ref{theorem2} hold. Then, for any sufficiently small
$\varepsilon_{0}>0$, there exists a constant $C>0$ such that, with probability
tending to one,
\[
Q_{n}(x) \leq C\,x^{-1/\gamma_{1}+\varepsilon_{0}}, \qquad\text{uniformly for
}1\leq x\leq T_{n,a}.
\]
Consequently,
\[
x^{-1}\left(  Q_{n}(x)\right)  ^{\alpha_{n}} \leq C\,x^{-1-\beta
/\gamma+\varepsilon_{0}}, \qquad\text{uniformly for }1\leq x\leq T_{n,a},
\]
with probability tending to one.
\end{lemma}

\begin{proof}
Recall that
\[
Q_{n}(x)=\frac{\overline{F}_{n}^{(\mathrm{NA})}(Z_{k}x)}{\overline{F}%
_{n}^{(\mathrm{NA})}(Z_{k})}.
\]
By Proposition~\ref{prop:app-stochastic} (iii), uniformly for $1\leq x\leq
T_{n,a}$,
\[
Q_{n}(x)=x^{-p/\gamma}\left[  1+O_{\mathbf{P}}\!\left(  k^{-1/2}%
x^{2\eta/\gamma+\varepsilon}\right)  +o_{\mathbf{P}}(1)\right]  .
\]
Since $p/\gamma=1/\gamma_{1}$, we obtain
\[
Q_{n}(x)=x^{-1/\gamma_{1}}\left[  1+O_{\mathbf{P}}\!\left(  k^{-1/2}%
x^{2\eta/\gamma+\varepsilon}\right)  +o_{\mathbf{P}}(1)\right]  .
\]
We now control the stochastic factor. By the truncation condition $x\leq
T_{n,a}$ and the assumptions on $a_{n}$, we have
\[
\sup_{1\leq x\leq T_{n,a}}k^{-1/2}x^{2\eta/\gamma+\varepsilon}=o_{\mathbf{P}%
}(1).
\]
Hence,
\[
\sup_{1\leq x\leq T_{n,a}}\left\vert 1+O_{\mathbf{P}}\!\left(  k^{-1/2}%
x^{2\eta/\gamma+\varepsilon}\right)  +o_{\mathbf{P}}(1)\right\vert
=O_{\mathbf{P}}(1).
\]
Therefore, for any sufficiently small $\varepsilon_{0}>0$, there exists a
constant $C>0$ such that, with probability tending to one,
\[
Q_{n}(x)\leq Cx^{-1/\gamma_{1}+\varepsilon_{0}},\qquad1\leq x\leq T_{n,a}.
\]
Next, since
\[
\alpha_{n}=\frac{\beta}{\widehat{p}_{k}}\quad\text{and}\quad\widehat{p}%
_{k}\overset{\mathbf{P}}{\longrightarrow}p,
\]
we have
\[
\alpha_{n}\leq\frac{\beta}{p}+\varepsilon_{0}%
\]
with probability tending to one. Consequently,
\[
\left(  Q_{n}(x)\right)  ^{\alpha_{n}}\leq Cx^{-\alpha_{n}/\gamma_{1}%
+\alpha_{n}\varepsilon_{0}}.
\]
Using again $1/\gamma_{1}=p/\gamma$ and the convergence of $\alpha_{n}$, we
obtain, after possibly reducing $\varepsilon_{0}$,
\[
-\frac{\alpha_{n}}{\gamma_{1}}+\alpha_{n}\varepsilon_{0}\leq-\frac{\beta
}{\gamma}+\varepsilon_{0},
\]
with probability tending to one. Hence,
\[
\left(  Q_{n}(x)\right)  ^{\alpha_{n}}\leq Cx^{-\beta/\gamma+\varepsilon_{0}%
},\qquad1\leq x\leq T_{n,a}.
\]
Multiplying by $x^{-1}$ yields
\[
x^{-1}\left(  Q_{n}(x)\right)  ^{\alpha_{n}}\leq Cx^{-1-\beta/\gamma
+\varepsilon_{0}},
\]
which completes the proof.\newpage
\end{proof}

\section{\textbf{Appendix B \label{Appendix B}}}%

%TCIMACRO{\FRAME{fhFU}{5.6593in}{2.7025in}{0pt}{\Qcb{Bias (left panels) and MSE
%(right panels) of $\protect\widehat{\gamma}_{1,k}^{\left(  NA\right)  }\left(
%1.01\right)  $ (blue line), $\protect\widehat{\gamma}_{1,k}^{\left(
%NA\right)  }\left(  1.5\right)  $ (green line), $\protect\widehat{\gamma
%}_{1,k}^{\left(  NA\right)  }\left(  2\right)  $ (red line),
%$\protect\widehat{\gamma}_{1,k}^{\left(  MNS\right)  }$ (purple line) and
%$\protect\widehat{\gamma}_{1,k}^{\left(  EFG\right)  }$ (black line) based on
%$2000$ samples of size $1000$ from the Burr model censored by the Burr
%distribution for $\gamma_{1}=0.4$ (top) and $\gamma_{1}=0.7$ (bottom), with
%$p=0.30.$ }}{\Qlb{fig1}}{bb-a1-04-07-p03.eps}%
%{\special{ language "Scientific Word";  type "GRAPHIC";  display "USEDEF";
%valid_file "F";  width 5.6593in;  height 2.7025in;  depth 0pt;
%original-width 16.6693in;  original-height 8.118in;  cropleft "0";
%croptop "1";  cropright "1";  cropbottom "0";
%filename 'bb-a1-04-07-p03.eps';file-properties "XNPEU";}} }%
%BeginExpansion
\begin{figure}[h]%
\centering
\includegraphics[
height=2.7025in,
width=5.6593in
]%
{bb-a1-04-07-p03.eps}%
\caption{Bias (left panels) and MSE (right panels) of $\protect\widehat{\gamma
}_{1,k}^{\left(  NA\right)  }\left(  1.01\right)  $ (blue line),
$\protect\widehat{\gamma}_{1,k}^{\left(  NA\right)  }\left(  1.5\right)  $
(green line), $\protect\widehat{\gamma}_{1,k}^{\left(  NA\right)  }\left(
2\right)  $ (red line), $\protect\widehat{\gamma}_{1,k}^{\left(  MNS\right)
}$ (purple line) and $\protect\widehat{\gamma}_{1,k}^{\left(  EFG\right)  }$
(black line) based on $2000$ samples of size $1000$ from the Burr model
censored by the Burr distribution for $\gamma_{1}=0.4$ (top) and $\gamma
_{1}=0.7$ (bottom), with $p=0.30.$ }%
\label{fig1}%
\end{figure}
%EndExpansion
%

%TCIMACRO{\FRAME{fhFU}{5.5971in}{2.7129in}{0pt}{\Qcb{Bias (left panels) and MSE
%(right panels) of $\protect\widehat{\gamma}_{1,k}^{\left(  NA\right)  }\left(
%1.01\right)  $ (blue line), $\protect\widehat{\gamma}_{1,k}^{\left(
%NA\right)  }\left(  1.5\right)  $ (green line), $\protect\widehat{\gamma
%}_{1,k}^{\left(  NA\right)  }\left(  2\right)  $ (red line),
%$\protect\widehat{\gamma}_{1,k}^{\left(  MNS\right)  }$ (purple line) and
%$\protect\widehat{\gamma}_{1,k}^{\left(  EFG\right)  }$ (black line) based on
%$2000$ samples of size $1000$ from the Burr model censored by the Burr
%distribution for $\gamma_{1}=0.4$ (top) and $\gamma_{1}=0.7$ (bottom) with
%$p=0.50.$}}{\Qlb{fig2}}{bb-a1-04-07-p05.eps}%
%{\special{ language "Scientific Word";  type "GRAPHIC";  display "USEDEF";
%valid_file "F";  width 5.5971in;  height 2.7129in;  depth 0pt;
%original-width 16.6693in;  original-height 8.118in;  cropleft "0";
%croptop "1";  cropright "1";  cropbottom "0";
%filename 'bb-a1-04-07-p05.eps';file-properties "XNPEU";}} }%
%BeginExpansion
\begin{figure}[h]%
\centering
\includegraphics[
height=2.7129in,
width=5.5971in
]%
{bb-a1-04-07-p05.eps}%
\caption{Bias (left panels) and MSE (right panels) of $\protect\widehat{\gamma
}_{1,k}^{\left(  NA\right)  }\left(  1.01\right)  $ (blue line),
$\protect\widehat{\gamma}_{1,k}^{\left(  NA\right)  }\left(  1.5\right)  $
(green line), $\protect\widehat{\gamma}_{1,k}^{\left(  NA\right)  }\left(
2\right)  $ (red line), $\protect\widehat{\gamma}_{1,k}^{\left(  MNS\right)
}$ (purple line) and $\protect\widehat{\gamma}_{1,k}^{\left(  EFG\right)  }$
(black line) based on $2000$ samples of size $1000$ from the Burr model
censored by the Burr distribution for $\gamma_{1}=0.4$ (top) and $\gamma
_{1}=0.7$ (bottom) with $p=0.50.$}%
\label{fig2}%
\end{figure}
%EndExpansion
%

%TCIMACRO{\FRAME{fhFU}{5.5971in}{2.6714in}{0pt}{\Qcb{Bias (left panel) and MSE
%(right panel) of $\protect\widehat{\gamma}_{1,k}^{\left(  NA\right)  }\left(
%1.01\right)  $ (blue line), $\protect\widehat{\gamma}_{1,k}^{\left(
%NA\right)  }\left(  1.5\right)  $ (green line), $\protect\widehat{\gamma
%}_{1,k}^{\left(  NA\right)  }\left(  2\right)  $ (red line),
%$\protect\widehat{\gamma}_{1,k}^{\left(  MNS\right)  }$ (purple line) and
%$\protect\widehat{\gamma}_{1,k}^{\left(  EFG\right)  }$ (black line) based on
%$2000$ samples of size $1000$ from Burr model censored by Burr for $\gamma
%_{1}=0.4$ (top) and $\gamma_{1}=0.7$ (bottom), with $p=0.70.$}}{\Qlb{fig3}%
%}{bb-a1-04-07-p07.eps}{\special{ language "Scientific Word";  type "GRAPHIC";
%display "USEDEF";  valid_file "F";  width 5.5971in;  height 2.6714in;
%depth 0pt;  original-width 16.6693in;  original-height 8.118in;
%cropleft "0";  croptop "1";  cropright "1";  cropbottom "0";
%filename 'bb-a1-04-07-p07.eps';file-properties "XNPEU";}} }%
%BeginExpansion
\begin{figure}[h]%
\centering
\includegraphics[
height=2.6714in,
width=5.5971in
]%
{bb-a1-04-07-p07.eps}%
\caption{Bias (left panel) and MSE (right panel) of $\protect\widehat{\gamma
}_{1,k}^{\left(  NA\right)  }\left(  1.01\right)  $ (blue line),
$\protect\widehat{\gamma}_{1,k}^{\left(  NA\right)  }\left(  1.5\right)  $
(green line), $\protect\widehat{\gamma}_{1,k}^{\left(  NA\right)  }\left(
2\right)  $ (red line), $\protect\widehat{\gamma}_{1,k}^{\left(  MNS\right)
}$ (purple line) and $\protect\widehat{\gamma}_{1,k}^{\left(  EFG\right)  }$
(black line) based on $2000$ samples of size $1000$ from Burr model censored
by Burr for $\gamma_{1}=0.4$ (top) and $\gamma_{1}=0.7$ (bottom), with
$p=0.70.$}%
\label{fig3}%
\end{figure}
%EndExpansion
%

%TCIMACRO{\FRAME{fhFU}{5.5763in}{2.7025in}{0pt}{\Qcb{Bias (left panels) and MSE
%(right panels) of $\protect\widehat{\gamma}_{1,k}^{\left(  NA\right)  }\left(
%1.01\right)  $ (blue line), $\protect\widehat{\gamma}_{1,k}^{\left(
%NA\right)  }\left(  1.5\right)  $ (green line), $\protect\widehat{\gamma
%}_{1,k}^{\left(  NA\right)  }\left(  2\right)  $ (red line),
%$\protect\widehat{\gamma}_{1,k}^{\left(  MNS\right)  }$ (purple line) and
%$\protect\widehat{\gamma}_{1,k}^{\left(  EFG\right)  }$ (black line) based on
%$2000$ samples of size $1000$ from the Fr\'{e}chet model censored by the
%Fr\'{e}chet distribution for $\gamma_{1}=0.4$ (top) and $\gamma_{1}=0.7$
%(bottom), with $p=0.30.$}}{\Qlb{fig4}}{ff-a1-04-07-p03.eps}%
%{\special{ language "Scientific Word";  type "GRAPHIC";  display "USEDEF";
%valid_file "F";  width 5.5763in;  height 2.7025in;  depth 0pt;
%original-width 16.6693in;  original-height 8.118in;  cropleft "0";
%croptop "1";  cropright "1";  cropbottom "0";
%filename 'ff-a1-04-07-p03.eps';file-properties "XNPEU";}} }%
%BeginExpansion
\begin{figure}[h]%
\centering
\includegraphics[
height=2.7025in,
width=5.5763in
]%
{ff-a1-04-07-p03.eps}%
\caption{Bias (left panels) and MSE (right panels) of $\protect\widehat{\gamma
}_{1,k}^{\left(  NA\right)  }\left(  1.01\right)  $ (blue line),
$\protect\widehat{\gamma}_{1,k}^{\left(  NA\right)  }\left(  1.5\right)  $
(green line), $\protect\widehat{\gamma}_{1,k}^{\left(  NA\right)  }\left(
2\right)  $ (red line), $\protect\widehat{\gamma}_{1,k}^{\left(  MNS\right)
}$ (purple line) and $\protect\widehat{\gamma}_{1,k}^{\left(  EFG\right)  }$
(black line) based on $2000$ samples of size $1000$ from the Fr\'{e}chet model
censored by the Fr\'{e}chet distribution for $\gamma_{1}=0.4$ (top) and
$\gamma_{1}=0.7$ (bottom), with $p=0.30.$}%
\label{fig4}%
\end{figure}
%EndExpansion
%

%TCIMACRO{\FRAME{fhFU}{5.5763in}{2.7233in}{0pt}{\Qcb{Bias (left panels) and MSE
%(right panels) of $\protect\widehat{\gamma}_{1,k}^{\left(  NA\right)  }\left(
%1.01\right)  $ (blue line), $\protect\widehat{\gamma}_{1,k}^{\left(
%NA\right)  }\left(  1.5\right)  $ (green line), $\protect\widehat{\gamma
%}_{1,k}^{\left(  NA\right)  }\left(  2\right)  $ (red line),
%$\protect\widehat{\gamma}_{1,k}^{\left(  MNS\right)  }$ (purple line) and
%$\protect\widehat{\gamma}_{1,k}^{\left(  EFG\right)  }$ (black line) based on
%$2000$ samples of size $1000$ from the Fr\'{e}chet model censored by the
%Fr\'{e}chet distribution for $\gamma_{1}=0.4$ (top) and $\gamma_{1}=0.7$
%(bottom), with $p=0.50.$}}{\Qlb{fig5}}{ff-a1-04-07-p05.eps}%
%{\special{ language "Scientific Word";  type "GRAPHIC";  display "USEDEF";
%valid_file "F";  width 5.5763in;  height 2.7233in;  depth 0pt;
%original-width 16.6693in;  original-height 8.118in;  cropleft "0";
%croptop "1";  cropright "1";  cropbottom "0";
%filename 'ff-a1-04-07-p05.eps';file-properties "XNPEU";}} }%
%BeginExpansion
\begin{figure}[h]%
\centering
\includegraphics[
height=2.7233in,
width=5.5763in
]%
{ff-a1-04-07-p05.eps}%
\caption{Bias (left panels) and MSE (right panels) of $\protect\widehat{\gamma
}_{1,k}^{\left(  NA\right)  }\left(  1.01\right)  $ (blue line),
$\protect\widehat{\gamma}_{1,k}^{\left(  NA\right)  }\left(  1.5\right)  $
(green line), $\protect\widehat{\gamma}_{1,k}^{\left(  NA\right)  }\left(
2\right)  $ (red line), $\protect\widehat{\gamma}_{1,k}^{\left(  MNS\right)
}$ (purple line) and $\protect\widehat{\gamma}_{1,k}^{\left(  EFG\right)  }$
(black line) based on $2000$ samples of size $1000$ from the Fr\'{e}chet model
censored by the Fr\'{e}chet distribution for $\gamma_{1}=0.4$ (top) and
$\gamma_{1}=0.7$ (bottom), with $p=0.50.$}%
\label{fig5}%
\end{figure}
%EndExpansion
%

%TCIMACRO{\FRAME{fhFU}{5.4829in}{2.7129in}{0pt}{\Qcb{Bias (left panels) and MSE
%(right panels) of $\protect\widehat{\gamma}_{1,k}^{\left(  NA\right)  }\left(
%1.01\right)  $ (blue line), $\protect\widehat{\gamma}_{1,k}^{\left(
%NA\right)  }\left(  1.5\right)  $ (green line), $\protect\widehat{\gamma
%}_{1,k}^{\left(  NA\right)  }\left(  2\right)  $ (red line),
%$\protect\widehat{\gamma}_{1,k}^{\left(  MNS\right)  }$ (purple line) and
%$\protect\widehat{\gamma}_{1,k}^{\left(  EFG\right)  }$ (black line) based on
%$2000$ samples of size $1000$ from the Fr\'{e}chet model censored by the
%Fr\'{e}chet distribution for $\gamma_{1}=0.4$ (top) and $\gamma_{1}=0.7$
%(bottom), with $p=0.70.$}}{\Qlb{fig6}}{ff-a1-04-07-p07.eps}%
%{\special{ language "Scientific Word";  type "GRAPHIC";  display "USEDEF";
%valid_file "F";  width 5.4829in;  height 2.7129in;  depth 0pt;
%original-width 16.6693in;  original-height 8.118in;  cropleft "0";
%croptop "1";  cropright "1";  cropbottom "0";
%filename 'ff-a1-04-07-p07.eps';file-properties "XNPEU";}} }%
%BeginExpansion
\begin{figure}[h]%
\centering
\includegraphics[
height=2.7129in,
width=5.4829in
]%
{ff-a1-04-07-p07.eps}%
\caption{Bias (left panels) and MSE (right panels) of $\protect\widehat{\gamma
}_{1,k}^{\left(  NA\right)  }\left(  1.01\right)  $ (blue line),
$\protect\widehat{\gamma}_{1,k}^{\left(  NA\right)  }\left(  1.5\right)  $
(green line), $\protect\widehat{\gamma}_{1,k}^{\left(  NA\right)  }\left(
2\right)  $ (red line), $\protect\widehat{\gamma}_{1,k}^{\left(  MNS\right)
}$ (purple line) and $\protect\widehat{\gamma}_{1,k}^{\left(  EFG\right)  }$
(black line) based on $2000$ samples of size $1000$ from the Fr\'{e}chet model
censored by the Fr\'{e}chet distribution for $\gamma_{1}=0.4$ (top) and
$\gamma_{1}=0.7$ (bottom), with $p=0.70.$}%
\label{fig6}%
\end{figure}
%EndExpansion
%

%TCIMACRO{\FRAME{fhFU}{5.5011in}{3.0926in}{0pt}{\Qcb{Bias (left panels) and MSE
%(right panels) of $\protect\widehat{\gamma}_{1,k}^{\left(  NA\right)  }\left(
%1.01\right)  $ (blue line), $\protect\widehat{\gamma}_{1,k}^{\left(
%NA\right)  }\left(  1.5\right)  $ (green line), $\protect\widehat{\gamma
%}_{1,k}^{\left(  NA\right)  }\left(  2\right)  $ (red line),
%$\protect\widehat{\gamma}_{1,k}^{\left(  MNS\right)  }$ (purple line) and
%$\protect\widehat{\gamma}_{1,k}^{\left(  EFG\right)  }$ (black line) based on
%$2000$ samples of size $1000$ from the Log-gamma model censored by the
%Log-gamma distribution for $\gamma_{1}=0.4$ (top) and $\gamma_{1}=0.7$
%(bottom), with $p=0.30.$}}{\Qlb{fig7}}{lg-a1-04-07-p03-g2.eps}%
%{\special{ language "Scientific Word";  type "GRAPHIC";  display "USEDEF";
%valid_file "F";  width 5.5011in;  height 3.0926in;  depth 0pt;
%original-width 10.6296in;  original-height 7.0843in;  cropleft "0";
%croptop "1";  cropright "1";  cropbottom "0";
%filename 'lg-a1-04-07-p03-g2.eps';file-properties "XNPEU";}} }%
%BeginExpansion
\begin{figure}[h]%
\centering
\includegraphics[
height=3.0926in,
width=5.5011in
]%
{lg-a1-04-07-p03-g2.eps}%
\caption{Bias (left panels) and MSE (right panels) of $\protect\widehat{\gamma
}_{1,k}^{\left(  NA\right)  }\left(  1.01\right)  $ (blue line),
$\protect\widehat{\gamma}_{1,k}^{\left(  NA\right)  }\left(  1.5\right)  $
(green line), $\protect\widehat{\gamma}_{1,k}^{\left(  NA\right)  }\left(
2\right)  $ (red line), $\protect\widehat{\gamma}_{1,k}^{\left(  MNS\right)
}$ (purple line) and $\protect\widehat{\gamma}_{1,k}^{\left(  EFG\right)  }$
(black line) based on $2000$ samples of size $1000$ from the Log-gamma model
censored by the Log-gamma distribution for $\gamma_{1}=0.4$ (top) and
$\gamma_{1}=0.7$ (bottom), with $p=0.30.$}%
\label{fig7}%
\end{figure}
%EndExpansion
%

%TCIMACRO{\FRAME{fhFU}{5.5089in}{3.0926in}{0pt}{\Qcb{Bias (left panels) and MSE
%(right panels) of $\protect\widehat{\gamma}_{1,k}^{\left(  NA\right)  }\left(
%1.01\right)  $ (blue line), $\protect\widehat{\gamma}_{1,k}^{\left(
%NA\right)  }\left(  1.5\right)  $ (green line), $\protect\widehat{\gamma
%}_{1,k}^{\left(  NA\right)  }\left(  2\right)  $ (red line),
%$\protect\widehat{\gamma}_{1,k}^{\left(  MNS\right)  }$ (purple line) and
%$\protect\widehat{\gamma}_{1,k}^{\left(  EFG\right)  }$ (black line) based on
%$2000$ samples of size $1000$ from the Log-gamma model censored by the
%Log-gamma distribution for $\gamma_{1}=0.4$ (top) and $\gamma_{1}=0.7$
%(bottom), with $p=0.50.$}}{\Qlb{fig8}}{lg-a1-04-07-p05-g2.eps}%
%{\special{ language "Scientific Word";  type "GRAPHIC";  display "USEDEF";
%valid_file "F";  width 5.5089in;  height 3.0926in;  depth 0pt;
%original-width 10.6296in;  original-height 7.0843in;  cropleft "0";
%croptop "1";  cropright "1";  cropbottom "0";
%filename 'lg-a1-04-07-p05-g2.eps';file-properties "XNPEU";}} }%
%BeginExpansion
\begin{figure}[h]%
\centering
\includegraphics[
height=3.0926in,
width=5.5089in
]%
{lg-a1-04-07-p05-g2.eps}%
\caption{Bias (left panels) and MSE (right panels) of $\protect\widehat{\gamma
}_{1,k}^{\left(  NA\right)  }\left(  1.01\right)  $ (blue line),
$\protect\widehat{\gamma}_{1,k}^{\left(  NA\right)  }\left(  1.5\right)  $
(green line), $\protect\widehat{\gamma}_{1,k}^{\left(  NA\right)  }\left(
2\right)  $ (red line), $\protect\widehat{\gamma}_{1,k}^{\left(  MNS\right)
}$ (purple line) and $\protect\widehat{\gamma}_{1,k}^{\left(  EFG\right)  }$
(black line) based on $2000$ samples of size $1000$ from the Log-gamma model
censored by the Log-gamma distribution for $\gamma_{1}=0.4$ (top) and
$\gamma_{1}=0.7$ (bottom), with $p=0.50.$}%
\label{fig8}%
\end{figure}
%EndExpansion
%

%TCIMACRO{\FRAME{fhFU}{5.5011in}{3.0926in}{0pt}{\Qcb{Bias (left panels) and MSE
%(right panels) of $\protect\widehat{\gamma}_{1,k}^{\left(  NA\right)  }\left(
%1.01\right)  $ (blue line), $\protect\widehat{\gamma}_{1,k}^{\left(
%NA\right)  }\left(  1.5\right)  $ (green line), $\protect\widehat{\gamma
%}_{1,k}^{\left(  NA\right)  }\left(  2\right)  $ (red line),
%$\protect\widehat{\gamma}_{1,k}^{\left(  MNS\right)  }$ (purple line) and
%$\protect\widehat{\gamma}_{1,k}^{\left(  EFG\right)  }$ (black line) based on
%$2000$ samples of size $1000$ from the Log-gamma model censored by the
%Log-gamma distribution for $\gamma_{1}=0.4$ (top) and $\gamma_{1}=0.7$
%(bottom), with $p=0.70.$}}{\Qlb{fig9}}{lg-a1-04-07-p07-g2.eps}%
%{\special{ language "Scientific Word";  type "GRAPHIC";  display "USEDEF";
%valid_file "F";  width 5.5011in;  height 3.0926in;  depth 0pt;
%original-width 10.6296in;  original-height 7.0843in;  cropleft "0";
%croptop "1";  cropright "1";  cropbottom "0";
%filename 'lg-a1-04-07-p07-g2.eps';file-properties "XNPEU";}} }%
%BeginExpansion
\begin{figure}[h]%
\centering
\includegraphics[
height=3.0926in,
width=5.5011in
]%
{lg-a1-04-07-p07-g2.eps}%
\caption{Bias (left panels) and MSE (right panels) of $\protect\widehat{\gamma
}_{1,k}^{\left(  NA\right)  }\left(  1.01\right)  $ (blue line),
$\protect\widehat{\gamma}_{1,k}^{\left(  NA\right)  }\left(  1.5\right)  $
(green line), $\protect\widehat{\gamma}_{1,k}^{\left(  NA\right)  }\left(
2\right)  $ (red line), $\protect\widehat{\gamma}_{1,k}^{\left(  MNS\right)
}$ (purple line) and $\protect\widehat{\gamma}_{1,k}^{\left(  EFG\right)  }$
(black line) based on $2000$ samples of size $1000$ from the Log-gamma model
censored by the Log-gamma distribution for $\gamma_{1}=0.4$ (top) and
$\gamma_{1}=0.7$ (bottom), with $p=0.70.$}%
\label{fig9}%
\end{figure}
%EndExpansion



\end{document}
